% Hebrew translation of the English 110 OBSERVER edition.
% Source: Dnarts/rule110-dependency-fronts, v2.2.1-monograph-multilingual.
% Build twice with XeLaTeX. Fonts: DejaVu Sans, Noto Serif, DejaVu Sans Mono.
\documentclass[11pt,a4paper]{article}
\usepackage[margin=24mm]{geometry}
\usepackage{fontspec}
\usepackage{amsmath,amssymb,amsthm}
\usepackage[english,bidi=default]{babel}
\babelprovide[import,main]{hebrew}
\babelfont{rm}{Noto Serif}
\babelfont[hebrew]{rm}{DejaVu Sans}
\babelfont{tt}[Scale=.87]{DejaVu Sans Mono}
\usepackage{array,booktabs,tabularx,enumitem}
\usepackage[table]{xcolor}
\definecolor{ink}{HTML}{1E293B}
\definecolor{slate}{HTML}{475569}
\definecolor{rulegray}{HTML}{CBD5E1}
\definecolor{tint}{HTML}{F8FAFC}
\definecolor{accent}{HTML}{174B69}
\usepackage[unicode,colorlinks=true,linkcolor=accent,urlcolor=accent,citecolor=accent,bookmarksnumbered=true]{hyperref}
\hypersetup{pdfauthor={Aleksandr Strand},pdflang={he},pdftitle={110 OBSERVER - יכולת התצפית על עברה של מערכת}}
\DeclareRobustCommand{\textenglish}[1]{\foreignlanguage{english}{#1}}
\newenvironment{english}{\begin{otherlanguage}{english}}{\end{otherlanguage}}
\addto\captionshebrew{%
 \renewcommand{\contentsname}{תוכן העניינים}%
 \renewcommand{\abstractname}{תקציר}%
 \renewcommand{\refname}{ביבליוגרפיה}%
 \renewcommand{\tablename}{טבלה}%
 \renewcommand{\proofname}{הוכחה}%
 \renewcommand{\appendixname}{נספח}%
}
\makeatletter
\def\tagform@#1{\maketag@@@{$\m@th(\mathrm{#1})$}}
\makeatother
\setlength{\emergencystretch}{2.5em}
\setcounter{tocdepth}{1}
\setlist{itemsep=1pt,topsep=3pt}
\newcommand{\src}[2]{\href{https://github.com/Dnarts/rule110-dependency-fronts/blob/v2.2.1-monograph-multilingual/#1}{#2}}
\newtheorem{theorem}{משפט}[section]
\newtheorem{lemma}[theorem]{למה}
\newtheorem{proposition}[theorem]{טענה}
\newtheorem{conjecture}[theorem]{השערה}
\theoremstyle{definition}
\newtheorem{definition}[theorem]{הגדרה}
\newtheorem{remark}[theorem]{הערה}
\begin{document}
\sloppy
\begin{center}
{\small\bfseries\color{accent}\textenglish{110 OBSERVER}}\\[7pt]
{\LARGE\bfseries יכולת התצפית על עברה של המערכת עצמה}\\[5pt]
{\large עלות המסר, עדויות סיבתיות וגבולות ההסקה}\\[10pt]
\textenglish{Aleksandr Strand} \hspace{1em} $\cdot$ \hspace{1em} 20 בנובמבר 2026
\end{center}
\noindent\textbf{תקציר.} כמה מידע על העבר דרוש למערכת להחלטתה הבאה, וכיצד נוכל לבדוק אם השתמשה בו? אנו נותנים קריטריון לפענוח מדויק וקובעים את עלות המסר. לאחר שתצורות הסריג הראשי של כלל 110 מתלכדות, צופה מחזיר קריאת חיישן שנשמרה באוגר חיצוני בן סיבית אחת. בעדכון הבא הסריג משחזר את הסיבית הקודמת שלו. ב־\textenglish{OLMo} עם משקולות קפואות, לפני שמוצגת מחרוזת חדשה, השולח מקודד את בחירתו הקודמת ואת המשוב בטוקן אחד המייצג כלל. ארבעת הצירופים של בחירה ומשוב מפיקים ארבעה מסרים, המשמשים שוב עבור מחרוזות חדשות. במחרוזות שסיביות הקצה שלהן שונות, המקבל פותר 64 מתוך 64 תנאי פיתוח באמצעות הטוקן שהועבר, לעומת 32 מתוך 64 עם הטוקן הקבוע \textenglish{A}. במערך הפיתוח המלא הוא עונה נכונה על 124 מתוך 128 תנאים ופותר 6 מתוך 8 בלוקים בשלמותם. התערבות בערוץ בוחנת את העברת הכלל. הסקה לגבי חוויה סובייקטיבית דורשת הנחות נוספות.

\begin{quote}\small
\textbf{היקף מהדורה זו.} המשפט בדבר הגאומטריה המדויקת של כלל 110 על הסריג האינסופי (\textenglish{\texttt{Rails.essential\_iff}}) הוכח ב־\textenglish{Lean 4.35.0}. בהכנת המונוגרפיה נבדק בנפרד גם \textenglish{\texttt{Rule110Observer.lean}}, עם שני הייבואים הקנוניים שלו. היומן מתעד את כל 11 דוחות האקסיומות (נספח~\ref{app:lean}). הגרסה הדו־שלבית של \textenglish{OLMo} השיגה \textenglish{\texttt{PASS\_DEV}} על שמונה מחרוזות אמצע. מערכי הבדיקה שהושארו בצד, \textenglish{9--20, 21--32, 33--44}, וערוץ האקטיבציות \textenglish{G3} מסומנים \textenglish{\texttt{NOT\_RUN}}. הנחת הגישור מסומנת \textenglish{\texttt{DECLARED\_NOT\_DERIVED}}.
\end{quote}
\tableofcontents
\clearpage

\section{השאלה, שתי דוגמאות וסימונים}
\emph{לפני שקובעים מה מערכת מסוגלת לעשות, כדאי לשאול איזה מידע לא עמד לרשותה.}
מתי מערכת דינמית יכולה לשחזר עובדה על עברה שלה? כמה סיביות עליה לקבל, ואיזו התערבות יכולה להראות שהשתמשה בהן? אנו מגדירים אילו גדלים זמינים למקבל בזמן הצגת השאילתה. שתי בניות בכלל 110 ממחישות מדוע זה חשוב.

\paragraph{שני חיישנים מכסים את החלון אך מחמיצים את ההבדל.}
נמקם חיישנים ב־$S=(-1,1)$ ונקרא אותם בזמנים 0 ו־1. המשימה היא לחזות את ערכו של תא 0 בזמן 2. קונוסי התלות של החיישנים מכסים את מלוא חלון התלות שלו, $[-2,2]$. במצב ההתחלתי הראשון רק בתא $-1$ יש 1. במצב השני יש 1 בתאים $-1$ ו־$0$. בשני המקרים החיישנים מדווחים פעמיים על אותו זוג $(1,0)$. בזמן 2 ערכי תא המטרה הם 0 ו־1, בהתאמה. כל מפענח הפועל לפי רישום זה יטעה באחד משני המקרים (סעיף~\ref{sec:observer-coverage}).

\paragraph{סיבית אחת משיבה הבחנה שאבדה.}
גדל המטרה כעת הוא הערך הקודם $Y$ של תא 5. בכל אחד משני זוגות ההיסטוריות, היפוך של $Y$ בלבד משנה את הקריאה $Z$ בתא 3 בזמן 2. הצופה שומר סיבית זו. בזמן 4 שתי תצורות הסריג בכל זוג כבר זהות לכל אורך הישר. נכתוב את $Z$ בתא $-1$. ערך שכנו השמאלי הוא 1, ושכנו הימני מחזיק את סיבית ההקשר $K$. העדכון הבא נותן $Z\oplus K=Y$. החזרת $1-Z$ לאותו מצב סריג הופכת את התשובה בכל ארבע ההיסטוריות (סעיף~\ref{sec:observer-return}).

בבנייה זו הערוץ נושא קריאת חיישן גולמית. בניסוי הדו־שלבי ב־\textenglish{OLMo}, השולח מחשב תחילה את הכלל המתוקן $p=h\oplus e$ מתוך בחירתו הקודמת $h$ והמשוב $e$. לפרק החדש נכנס רק הטוקן $m\in\{\mathrm{A},\mathrm{B}\}$, שנבחר לפני הצגת המחרוזת $\mathbf{x}$. המקבל מעבד את הטוקן מחדש עם מטמון \textenglish{KV} ריק. כך נמדדת העברת כלל שכבר נבחר מעבר לגבול שבין פרקים. תיקון פנימי של האסטרטגיה כאשר המצב נשמר ברציפות דורש בדיקה נפרדת (סעיף~\ref{sec:olmo}).

\begin{center}\small
\begin{tabularx}{\linewidth}{@{}lX@{}}\toprule
סימון & משמעות במונוגרפיה \\
\midrule
$\omega,H,W,Z$ & היסטוריה קבילה, מידע המקבל, המידע הזמין לצופה והמסר שנמסר.\\
$Y$ & סיבית המטרה על העבר. בהצבה הכללית יותר של \textenglish{T1--T3}, האלפבית שלה יכול להיות כל קבוצה סופית.\\
$t,i,j,s$ & זמן בדיד וקואורדינטות שלמות בסריג. $e_j$ הופך רק את התא ההתחלתי $j$.\\
$B_t,\widetilde B_t,S_t$ & התצורה הבלתי מופרעת, התצורה המופרעת וקבוצת הקואורדינטות שבהן הן נבדלות.\\
$G_t,D_t,I_t$ & ערך התא בראשית, הקואורדינטות ההתחלתיות המהותיות וחלון התלות התאורטי.\\
$h,e,p,m$ & הבחירה הקודמת של \textenglish{OLMo}, סיבית המשוב, הכלל המתוקן וטוקן המסר.\\
$\mathbf{x}=LwR$ & קלט חדש באורך שמונה סיביות, ובו אמצע $w$ באורך שש סיביות וסיביות קצה $L,R$.\\
$B,E(a,t)$ & חבילת הנחות הגישור והפרדיקט של חוויה פנומנלית עבור המופע $a$ בזמן $t$.\\
\bottomrule\end{tabularx}
\end{center}

נוסחה קובעת מה נובע מהנחות מוגדרות (\textenglish{\texttt{COMPUTED}}). טענה על מכונה מסוימת דורשת עדות שניתן להריץ או דוח מדידה (\textenglish{\texttt{MEASURED}}). הסטטוס \textenglish{\texttt{UNBOUND}} מציין שחסר קישור נדרש לראיות. הסטטוס \textenglish{\texttt{NOT\_RUN}} מציין שלא בוצעה הרצה בגרסה המתוארת. בהמשך, ״מדויק״ מתייחס לכל איברי תחום מתמטי מוגדר או למערך בדיקות סופי המצוין במפורש. סעיף~\ref{sec:observer-theory} פותח בקביעת התנאים שבהם קיים מפענח מדויק.

\section{גישה לעבר: תצפיות, מסרים וזיכרון}
\label{sec:observer-theory}

תהי $U\ne\varnothing$ קבוצת ההיסטוריות הקבילות של התקן אחד בגבולות מערכת ובפרוטוקול קבועים. סיבית המטרה $Y(u)\in\{0,1\}$ מתארת את עברו של ההתקן בהיסטוריה $u$. נסמן ב־$H(u)$ את המידע הזמין לנמען בעת ההכרעה, ב־$W(u)$ את המידע הזמין לשולח וב־$Z(u)$ את המסר שנמסר. אפשר לכלול ב־$H$ קריאות קודמות רק אם הנמען יכול לגשת אליהן דרך הממשק שלו או אם שמר אותן בזיכרונו. שני פרוטוקולים המעניקים גישה שונה לתיעוד מגדירים משימות שונות.

\subsection{\textenglish{T1}. תנאי לפענוח מדויק}
\label{thm:observer-t1}

לכל העתקת תצפית $O:U\to\mathcal O$, שחזור דטרמיניסטי מדויק של $Y$ בכל ההיסטוריות הקבילות אפשרי אם ורק אם
\begin{equation}
 \forall u,v\in U\qquad O(u)=O(v)\ \Longrightarrow\ Y(u)=Y(v).
 \label{eq:observer-fibres}
\end{equation}
\emph{הוכחה.} אם $Y=d\circ O$, החלת אותה פונקציה $d$ על ארגומנטים שווים נותנת תשובות שוות. בכיוון ההפוך, לכל $o\in O(U)$ הקבוצה $\{Y(u):O(u)=o\}$ אינה ריקה, ולפי ההנחה יש בה איבר יחיד. נגדיר את $d(o)$ להיות אותו איבר. ערכיה של $d$ מחוץ ל־$O(U)$ שרירותיים. בכך הוכחנו שקיימת פונקציה כזאת. כאשר התחום סופי, אפשר לממש אותה כטבלת חיפוש, אם יש די זיכרון ויכולת מיעון.

התנאי חל בנפרד על $O=H$ ועל $O=(H,Z)$. כדי להראות שקבלת מסר משפרת את הגישה ללא שגיאות, די להציג זוג היסטוריות שלהן אותו $H$ וערכי $Y$ שונים, ואז לתת מפענח מדויק מתוך $(H,Z)$ עבור \emph{כל} מחלקת ההיסטוריות. למשפחת תצפיות $O_t$, המועד הראשון שבו מתאפשרת גישה כזאת הוא ה־$t$ הקטן ביותר המקיים את \eqref{eq:observer-fibres}, אם הוא קיים. במערכת שיש בה מחיקה, השאלה אם התנאי מוסיף להתקיים בהמשך היא שאלה נפרדת.

\subsection{\textenglish{T2}. עלותו של מסר יחיד}
\label{thm:observer-t2}

נניח ש־$U$ סופית. השולח מעביר מילה אחת באורך קבוע $c(W)$ בערוץ ללא רעש. זמני השידור והמסירה נקבעים מראש. נבנה גרף לא מכוון $G$ שקבוצת קודקודיו היא $W(U)$. נחבר את $w,w'$ בקשת אם קיימים $u,v\in U$ כך שמתקיים
\[
 W(u)=w,\qquad W(v)=w',\qquad H(u)=H(v),\qquad Y(u)\ne Y(v).
\]
אנו מתירים $w=w'$. אם קיימת לולאה, שום מסר מן הצורה הזאת אינו מספיק. אם הגרף חסר לולאות, גודל האלפבית המזערי והמספר המזערי של ספרות בינאריות שיש לשדר הם, בהתאמה,
\[
 M_{\min}=\chi(G),\qquad b_{\min}=\lceil\log_2\chi(G)\rceil .
\]
\emph{הוכחה.} שידור אותו מסר בשני קצות קשת היה יוצר צמד $(H,Z)$ זהה עבור מטרות שונות, בסתירה ל־\textenglish{T1}. לכן ערכיו של כל קוד מדויק מגדירים צביעה נאותה. לולאה הייתה מחייבת לצבוע קודקוד אחד בשני צבעים שונים. בכיוון ההפוך, בכל צביעה נאותה שתי היסטוריות שלהן אותו $H$ ומטרות שונות מקבלות צבעים שונים. כשהשולח משדר את הצבע, \textenglish{T1} מבטיח שקיים מפענח מדויק. המספר המזערי של סיביות לקידוד באורך קבוע של $M$ צבעים הוא בדיוק $\lceil\log_2 M\rceil$, ובכלל זה המסר הריק כאשר $M=1$.

החסם התחתון $\max_{h\in H(U)}|\{Y(u):H(u)=h\}|$ ניתן להשגה אם השולח יודע גם את $H$ וגם את $Y$. הוא יכול למספר את המטרות האפשריות בנפרד לכל $h$. כאשר הגישה להיסטוריה מוגבלת, חסם זה עשוי להיות קטן ממש מ־$M_{\min}$. נניח שהנמען מחזיק באחד הזוגות הסדורים $ab,ac,bc$. השולח יודע איזו אות בזוג הזה נבחרה קודם, וסיבית המטרה מציינת אם נבחר המקום הראשון או השני. האותיות בכל זוג דורשות צבעים שונים. שלוש האותיות יוצרות משולש, ולכן נחוצים שלושה סמלי מסר אף שלסיבית המטרה יש רק שני ערכים אפשריים. הקשר בין קידוד חד־פעמי ללא שגיאות עם מידע צד לבין צביעת גרפים מקורו בעבודתו של \href{https://doi.org/10.1109/TIT.1976.1055607}{\textenglish{Witsenhausen}}. גם \href{https://doi.org/10.3390/e26040338}{\textenglish{Charpenay, Le Treust, and Roumy}} עוסקים בבעיית חישוב פונקציה בנוכחות מידע צד.

\subsection{\textenglish{T3}. שגיאות בערוץ ועלותה של שאילתה ממוענת}
\label{thm:observer-t3}

נניח שבהיסטוריה $u$ האות המשובש יכול לקבל כל ערך בקבוצה לא ריקה $\mathcal R(u)$. שחזור תחת \emph{כל} שיבוש קביל אפשרי בדיוק כאשר
\[
 H(u)=H(v),\quad Y(u)\ne Y(v)
 \quad\Longrightarrow\quad \mathcal R(u)\cap\mathcal R(v)=\varnothing .
\]
סמל השייך לחיתוך היה מחייב לתת שתי תשובות שונות לאותה תצפית. אם כל החיתוכים המצוינים ריקים, לכל צמד $(h,r)$ שניתן להגיע אליו יש ערך מטרה יחיד. זוהי הוכחת \textenglish{T1} כאשר מרחיבים את ההיסטוריות כך שיכללו את תוצאת הערוץ.

עבור התפלגות משותפת סופית של $(H,R,Y)$ תחת הפסד אפס־אחד, הסתברות השגיאה המזערית מבטאת תנאי אחר, והיא שווה ל־
\[
 P_e^*=1-\sum_{h,r}\max_{y\in\{0,1\}}\Pr(H=h,R=r,Y=y).
\]
ואכן, תשובה דטרמיניסטית $d(h,r)$ נותנת הסתברות הצלחה $\sum_{h,r}\Pr(H=h,R=r,Y=d(h,r))$. אפשר למקסם כל מחובר בנפרד. לדוגמה, נניח ש־$H,Z$ הן סיביות בלתי תלויות והוגנות, ש־$Y=H\oplus Z$, וש־$R=Z\oplus N$ עם רעש בלתי תלוי $N\sim\mathrm{Bernoulli}(p)$, $0\le p\le 1/2$. לתשובה המיטבית $H\oplus R$ הסתברות שגיאה $p$. כאשר $p>0$, אין עוד ערובה לשחזור מדויק תחת כל שיבוש.

יתרון האינטראקציה תלוי בסוגי המסרים המותרים. יהיו $A_0,A_1$ שתי סיביות שנרשמו קודם, ללא אילוצים על ערכיהן. כתובת $K$ הידועה לנמען מציינת את הסיבית הדרושה מן העבר, $Y=A_K$. כל שמונה השלשות $(K,A_0,A_1)$ קבילות. סיבית שאילתה אחת $K$ וסיבית תשובה אחת $A_K$ מספיקות. נניח במקום זאת שסיבית התשובה חייבת להיות פונקציה קבועה $c(A_0,A_1)$ שאינה תלויה בכתובת. בהינתן ערך אחד של הפונקציה הזאת, ניתן היה להפעיל מפענח מדויק עם $K=0$ ועם $K=1$ כדי לשחזר את הצמד $(A_0,A_1)$ כולו. ארבעת הצמדים השונים דורשים לפחות ארבעה מסרים. לפיכך שידור הכתובת מקטין את אלפבית התשובה הדרוש לפי פרוטוקול תקשורת זה. שתי סיביות תשובה כבר מספיקות לשדר את כל התיעוד בלי לשאול. עבור שאילתה קבועה יחידה, השמעה חוזרת של תשובה נכונה באותו תוכן ובאותו מועד מסירה יוצרת אותה תוצאה אצל נמען דטרמיניסטי.

\subsection{\textenglish{T4}. דיווח עצמי ושימור מצב}
\label{thm:observer-t4}

נבחן שני התקנים, ששניהם מחשבים במדויק אותה תשובה $d(H,Z)=Y$ ומפרסמים אותה. הראשון שומר סיבית $M=Y$. השני מוחק את העדות ומגיע לאותו מצב נגיש עבור לפחות שתי היסטוריות שבהן $Y$ שונה. בשלב זה נחסמת הגישה לדיווח הפומבי הקודם ולערוץ של הצופה. שני ההתקנים מקבלים סיבית חדשה $C$ ומתבקשים לבצע את הפעולה $Y\oplus C$. הראשון מחשב $M\oplus C$. עבור ההתקן השני, נקבע אותו $C$ בשתי ההיסטוריות שתוארו. כל הקלטים הזמינים זהים, אך הסיביות הנדרשות הפוכות. לפי \textenglish{T1}, להתקן השני אין מפענח מדויק. לפיכך הפעולה המושהית מבחינה בין שני מימושים שהפיקו דיווחים מיידיים זהים.

התערבות ב־$Z$ כאשר המצב הקודם נשאר קבוע בודקת אם הפעולה תלויה באות שנמסר. כדי לבדוק אם התוצאה עוברת לאפיזודה אחרת, נדרשת התערבות שנייה: מחליפים או מוחקים את $M$, משאירים את השאילתה החדשה קבועה וחוסמים את הגישה לדיווח הקודם. אלה שתי השוואות נגד־עובדתיות שונות. במשימה זו אוגר פנימי אחד מספיק לביצוע הפעולה המושהית הנכונה.

\section{הגאומטריה של הקואורדינטות המהותיות בכלל 110}\label{sec:geometry}
\subsection{הגודל הנצפה ועדויות באמצעות היפוך סיבית יחידה}
נקבע תא מטרה בראשית הצירים ואופק זמן $t$. קואורדינטת קלט היא מהותית אם קיימת תצורת התחלה שבה היפוך הסיבית באותה קואורדינטה בלבד משנה את הפלט. הרקע עשוי להיות תלוי הן בקואורדינטה והן באופק הזמן. סדר הכמתים הזה מבחין בין רגישות במקרה הגרוע ביותר לבין מהירותו של גליידר מסוים. בכלל~110 ידועות בניות של חישוב אוניברסלי \cite{cook2004} וכן סיווגים של מבנים נעים \cite{wolfram1986,martinez2006}. הקשר בין תצפיות לבין מחלקות חיזוי ומגבלות על ההיסק נדון ב־\cite{shalizi2001,wolpert2008}.

\begin{definition}[כלל 110 והגודל הנצפה]
יהי $X = \{0, 1\}^{\mathbb{Z}}$. הצורה האלגברית הנורמלית של הכלל המקומי מעל $\mathbb{F}_2$ היא
\begin{equation}\label{eq:rule110}
f(l, c, r) = c \oplus r \oplus cr \oplus lcr.
\end{equation}
העדכון הסינכרוני $F : X \to X$ נתון על ידי $(F(u))_i = f(u_{i-1}, u_i, u_{i+1})$. הגודל הנצפה בראשית הצירים בזמן $t \in \mathbb{N}$ הוא $G_t(u) = (F^t(u))_0$.
\end{definition}

\begin{definition}[פירוק לצורך חיזוי ומשתנים מהותיים]
מפת תצפית $O : X \to \mathcal{O}$ מספיקה עבור $G_t$ אם קיימת מפה $P : \mathcal{O} \to \{0, 1\}$ כך ש־$G_t = P \circ O$. עבור $j \in \mathbb{Z}$, נסמן את היפוכה של סיבית יחידה באמצעות $(u \oplus e_j)_i = u_i \oplus \mathbf{1}_{i=j}$. קבוצת הקואורדינטות ההתחלתיות המהותיות היא
\begin{equation}\label{eq:essential_set}
D_t = \{ j \in \mathbb{Z} : \exists u \in X,\; G_t(u) \ne G_t(u \oplus e_j) \}.
\end{equation}
\end{definition}

מכיוון ש־$G_t$ תלוי במספר סופי של קואורדינטות בלבד, $D_t$ היא קבוצת קואורדינטות הקלט היחידה המינימלית ביחס להכלה שמספיקה לחישוב $G_t$. כל קטע תצפית מספיק חייב להכיל את הקטע השלם הקטן ביותר המכיל את כל איברי $D_t$.

כדי לבנות עדויות על רקעים מחזוריים, אנו מבחינים בין קואורדינטת הקלט $j$ לבין הקואורדינטה הפיזית $s$ שבה נצפה פגם. עבור מילה מחזורית $w$, מופע מרחבי $\phi$ והתצורה האינסופית המתאימה $B_0^{w,\phi}$, נגדיר
\begin{equation}\label{eq:diff_support}
B_t^{w,\phi} = F^t(B_0^{w,\phi}), \qquad
\widetilde B_t^{w,\phi} = F^t(B_0^{w,\phi} \oplus e_0), \qquad
S_t^{w,\phi} = \{ s \in \mathbb{Z} : B_t^{w,\phi}(s) \ne \widetilde B_t^{w,\phi}(s) \}.
\end{equation}
העדכון מתחלף עם הזזות, ולכן $F \circ \sigma^d = \sigma^d \circ F$, כאשר $(\sigma^d u)_i = u_{i-d}$. לכל $s \in S_t^{w,\phi}$, הזזת זוג תצורות ההתחלה ב־$-s$ נותנת עדות לכך ש־$j = -s \in D_t$. מכאן שלכל משפחה $\mathcal{W}$ של רקעים ומופעים מתקיים
\begin{equation}\label{eq:coverage_bridge}
C_t(\mathcal{W}) = \bigcup_{(w,\phi) \in \mathcal{W}} \{ -s : s \in S_t^{w,\phi} \} \subseteq D_t.
\end{equation}
לפיכך, הקצה השמאלי $j = -L(t)$ של חלון הקלט מתאים לחזית הפיזית הימנית $s = L(t)$ של הפגם. הקצה הימני $j = t$ מתאים לחזית הפיזית השמאלית $s = -t$.

\subsection{חלון התלות המדויק על הסריג האינסופי}

נסמן $L(t) = t - \lfloor t/3 \rfloor$ ו־$I_t = [-L(t), t] \cap \mathbb{Z}$.

\begin{lemma}[חסם תלת־צעדי על התפשטות ימינה]\label{lem:three_step}
עבור מילה בת שבע סיביות $b = (b_0, \dots, b_6) \in \{0, 1\}^7$, יהי $Q(b)$ הערך בתא המרכזי שלה לאחר שלושה צעדים סינכרוניים של כלל 110. אזי
\begin{equation}\label{eq:seven_bit}
Q(b_0, b_1, \dots, b_6) = Q(1 \oplus b_0, b_1, \dots, b_6).
\end{equation}
\end{lemma}

\begin{proof}
לפי \eqref{eq:rule110}, הנגזרת הבוליאנית לפי הארגומנט השמאלי היא
\begin{equation}\label{eq:bool_deriv}
f(1, c, r) \oplus f(0, c, r) = cr.
\end{equation}
נסמן ב־$b_i(t)$ את ערך התא $i$ באחד ממסלולי ההתפתחות. נניח ששתי תצורות זהות בכל התאים שמימין ל־$p$, ושההבדל הימני ביותר ביניהן מתקדם תא אחד בכל אחד משלושה צעדים, לאורך $p \to p+1 \to p+2 \to p+3$. ההתקדמות הראשונה מחייבת שהערכים המשותפים בתאים $p+1$ ו־$p+2$ יהיו $1$ ב־$t=0$. ההתקדמות השנייה מחייבת שהערכים בתאים $p+2$ ו־$p+3$ יהיו $1$ ב־$t=1$. מכיוון ש־$f(1, 1, z) = 1 \oplus z$, בהכרח $b_{p+3}(0) = 0$. מן הזהות $f(1, 0, z) = z$ נובע אז שבהכרח $b_{p+4}(0) = 1$. לכן $b_{p+4}(1) = f(0, 1, b_{p+5}(0)) = 1$, ומתקבל $b_{p+3}(2) = f(1, 1, 1) = 0$. מכאן שההתקדמות השלישית אינה אפשרית. גם ליבת \textenglish{Lean 4} מאמתת את \eqref{eq:seven_bit} בבדיקת כל $2^7 = 128$ הקלטים.
\end{proof}

\begin{theorem}[חסם אוניברסלי על קבוצת התלות]\label{thm:upper_bound}
לכל $t \in \mathbb{N}$ מתקיים $D_t \subseteq I_t = [-(t - \lfloor t/3 \rfloor), t] \cap \mathbb{Z}$.
\end{theorem}

\begin{proof}
אם שתי תצורות זהות בכל התאים שמימין ל־$p_0$, לאחר צעד אחד הן עדיין זהות בכל התאים שמימין ל־$p_0 + 1$, ולפי הלמה~\ref{lem:three_step}, לאחר שלושה צעדים הן זהות בכל התאים שמימין ל־$p_0 + 2$. נכתוב $t = 3q + r$, כאשר $q = \lfloor t/3 \rfloor$ ו־$r \in \{0, 1, 2\}$. הזהות נשמרת בכל התאים שמימין ל־$p_0 + 2q + r = p_0 + L(t)$. לכן הבדל ב־$j < -L(t)$ אינו יכול להגיע לראשית הצירים עד הזמן $t$. עבור $j > t$, החסם נובע מחרוט האור הרגיל של כלל בעל רדיוס אחד.
\end{proof}

\begin{theorem}[חדות החסם בשני הקצוות]\label{thm:sharpness}
לכל $t \in \mathbb{N}$, שני הקצוות $-L(t)$ ו־$t$ שייכים ל־$D_t$. קטע התצפית המספיק הקטן ביותר הוא $I_t$. רוחבו מקיים
\begin{equation}
m_{\min}(t) = 2t + 1 - \lfloor t/3 \rfloor, \qquad \lim_{t \to \infty} \frac{m_{\min}(t)}{t} = \frac{5}{3}.
\end{equation}
\end{theorem}

\begin{proof}
עבור הקצה הימני $j = t$, נשווה את תצורת האפסים $0^{\mathbb{Z}}$ לתצורה שבה יש סיבית אחת שערכה אחד, במיקום $t$, ונסמן אותה ב־$e_t$. כיוון ש־$f(0, 0, 0) = 0$ ו־$f(0, 0, 1) = 1$, הסיבית השמאלית ביותר שערכה אחד נעה בדיוק תא אחד שמאלה בכל צעד ומגיעה לראשית הצירים לאחר $t$ צעדים.

עבור הקצה השמאלי $j = -L(t)$, ניקח את הרקע בעל המחזור ארבע, $B_0(i) = (1, 0, 1, 1)[i \bmod 4]$, ואת ההפרעה שלו $\widetilde B_0 = B_0 \oplus e_2$. טבלה~\ref{tab:period4_defect} מציגה את תומך ההפרש $S_t$, המתקבל מהפעלת כלל 110 במשך שישה צעדים.

\begin{table}[h]
\centering
\caption{פגם יחיד על הרקע $(1011)$. לאחר שלב המעבר ההתחלתי, זוג התצורות חוזר על עצמו בהעתק $(\Delta t, \Delta s) = (3, +2)$.}
\label{tab:period4_defect}
\begin{tabular}{cccc}
\toprule
$t$ & מילת הרקע $B_t$ (מחזור $4$) & תומך ההפרש $S_t$ & $\max S_t$ \\
\midrule
$0$ & $1011$ & $\{2\}$ & $2$ \\
$1$ & $1110$ & $\{1, 3\}$ & $3$ \\
$2$ & $1011$ & $\{1, 3, 4\}$ & $4$ \\
$3$ & $1110$ & $\{0, 1, 4\}$ & $4$ \\
$4$ & $1011$ & $\{-1, 0, 1, 3, 5\}$ & $5$ \\
$5$ & $1110$ & $\{3, 5, 6\}$ & $6$ \\
$6$ & $1011$ & $\{2, 3, 6\}$ & $6$ \\
\bottomrule
\end{tabular}
\end{table}

מן המקומיות נובע שהשורות ב־$t=3$ וב־$t=6$ זהות עד כדי הזזה של שני תאים ימינה, ומתקבלות הזהויות $B_6 = \sigma^2 B_3$ ו־$\widetilde B_6 = \sigma^2 \widetilde B_3$. העובדה שהעדכון מתחלף עם הזזות נותנת $(B_{3+3k+r}, \widetilde B_{3+3k+r}) = \sigma^{2k}(B_{3+r}, \widetilde B_{3+r})$ לכל $k \ge 0$ ו־$r \in \{0, 1, 2\}$. בצירוף השורות ההתחלתיות, מתקבל $\max S_t = 2 + L(t)$ לכל $t \ge 0$. הזזה של הזוג ההתחלתי ב־$-[2 + L(t)]$ מעבירה את היפוך הסיבית ההתחלתי אל $-L(t)$ ואת ההבדל בזמן $t$ אל ראשית הצירים. לכן $-L(t) \in D_t$.
\end{proof}

\subsection{טבעות סופיות, קיבולת קידוד מצבים וגישה לקלט}

על הטבעת הסופית $\mathbb{Z}_N$ ($N \ge 3$), יהי $X_N = \{0, 1\}^{\mathbb{Z}_N}$ מרחב התצורות ותהי $F_N : X_N \to X_N$ מפת העדכון. קבוצת המצבים שאליהם ניתן להגיע בזמן $t$ היא $\mathcal{R}_t^{(N)} = F_N^t(X_N)$, ומספר הסיביות הדרוש לקידוד מצב מתוך קבוצה זו הוא $C_t^{(N)} = \lceil \log_2 |\mathcal{R}_t^{(N)}| \rceil$. נגדיר את הדינמיקה על זוגות $\Phi_N(u, v) = (F_N(u), F_N(v))$, את קבוצת הזוגות ההתחלתיים $E_j = \{ (u, u \oplus e_j) : u \in X_N \}$, ואת קבוצת הזוגות שנבדלים בראשית הצירים, $W_0 = \{ (u,v) : u_0 \ne v_0 \}$. אזי
\begin{equation}
j \in D_t^{(N)} \iff \Phi_N^t(E_j) \cap W_0 \ne \emptyset.
\end{equation}

\begin{proposition}[תעודה סופית למהותיות מתמשכת]\label{prop:ring_cert}
יהי $\mu_N$ האורך המרבי של שלב המעבר במערכת הדטרמיניסטית הסופית $F_N$, ויהי $P_N$ הכפולה המשותפת הקטנה ביותר של אורכי המחזורים שלה. אזי $D_{t+P_N}^{(N)} = D_t^{(N)}$ לכל $t \ge \mu_N$. לפיכך, לכל $T \in \mathbb{N}$, התנאי $\forall t \ge T,\; D_t^{(N)} = \mathbb{Z}_N$ שקול לבדיקה ש־$\Phi_N^t(E_j) \cap W_0 \ne \emptyset$ לכל $j \in \mathbb{Z}_N$ ולכל $T \le t < \max(T, \mu_N) + P_N$.
\end{proposition}

\begin{proof}
עבור $t\ge\mu_N$, כל מצב שאליו ניתן להגיע, $F_N^t(u)$, נמצא על מחזור שאורכו מחלק את $P_N$. לכן המפות $F_N^{t+P_N}$ ו־$F_N^t$ זהות בכל $X_N$. לכל $j$, ובכל זוג $(u,u\oplus e_j)$, שוויון או אי־שוויון בין ערכי התא בראשית הצירים נשמר בשני הזמנים. החל מהזמן $\max(T,\mu_N)$ די לבדוק מחזור אחד. אם $T<\mu_N$, יש לבדוק גם כל זמן מ־$T$ עד $\mu_N-1$.
\end{proof}

\begin{remark}[דחיסת מצבים ורגישות לקלט]\label{rem:ring5}
רגישות על הסריג האינסופי אינה נשמרת בהכרח על טבעת סופית. בכלל 110, עבור $N = 5$, כל $32$ המצבים מגיעים ל־$0^5$ עד הזמן $t = 5$, ולכן $D_t^{(5)} = \emptyset$ לכל $t \ge 5$. על הטבעת שבה $N = 20$, מנייה ממצה ב־$t = 121$ נותנת $|\mathcal{R}_{121}^{(20)}| = 2095$ מצבים שאליהם ניתן להגיע, ולקידודם דרושות $12$ סיביות. באותו זמן מתקיים $D_{121}^{(20)} = \mathbb{Z}_{20}$, כך שכל $20$ סיביות הקלט מהותיות. דוגמאות אלה מבחינות בין ארבעה היבטים בלתי תלויים של חיזוי: \emph{סיבתיות} (התומך $D_t$), \emph{גישה} (קריאת קואורדינטות יחידות או שימוש בנתונים שעברו עיבוד מקדים), \emph{קיבולת} ($C_t^{(N)}$) ו\emph{סיבוכיות חישובית}.
\end{remark}

נחזור לסריג האינסופי. לכל קואורדינטה בפנים החלון, כמו גם בכל אחד מקצותיו, נוכל למצוא זוג תצורות התחלה המשמש עדות להבדל בתא המטרה.

\begin{theorem}[כל קואורדינטה בקטע היא מהותית]\label{thm:gap-free}
לכל $t \in \mathbb{N}$ מתקיים $D_t = I_t = [-(t - \lfloor t/3 \rfloor), t] \cap \mathbb{Z}$.
\end{theorem}

ההוכחה מובאת בנספח~\ref{app:lean}. עבור $t<2485$, תעודה סופית מספקת עדות לכל זוג $(t,j)$ בחלון. עבור $t\ge2485$, חמש רצועות סמוכות מכסות את החלון (סעיף~\ref{sec:five-bands}).

\paragraph{מסקנה לגבי הגישה לנתונים.}
כל קבוצה קבועה של קואורדינטות התחלתיות המספיקה לחישוב מדויק של $G_t$ על כל התצורות מכילה את $I_t$. המקומיות והחסם העליון מראים ש־$I_t$ עצמו מספיק. לכן המספר המזערי של קואורדינטות בתצפית כזאת הוא $2t+1-\lfloor t/3\rfloor$. העדות יכולה להיות תלויה הן ב־$t$ והן ב־$j$. טענות על עלותן של שאילתות אדפטיביות, על עיבוד מקדים או על זיכרון דורשות מודל חישובי נפרד.

\section{כיסוי ויכולת הבחנה: עדויות קונסטרוקטיביות בכלל~110}
\label{sec:observer-witnesses}

נבחן שתי שאלות על אותה דינמיקה. בשאלה הראשונה, צופה מקבל יומן קריאות מכמה אתרים וחוזה את ערכו העתידי של תא. בשנייה, המקבל משחזר את סיבית העבר שלו $Y$ לאחר שהיומן המקומי חדל להבחין בין ערכיה הקודמים. בשני המקרים אנו בוחנים את הקריטריון מסעיף~\ref{sec:observer-theory}: ערך היעד צריך להיות קבוע על כל סיב של העתקת התצפית.

\subsection{חיזוי מתוך יומן ומגבלות הכיסוי הגאומטרי}
\label{sec:observer-coverage}

לתצורה התחלתית $b$, נסמן ב־$b_t(s)$ את ערכה באתר $s$ בזמן $t$. תהי $S$ רשימה סופית של אתרים קבועים, ויהי $\tau$ זמן התצפית האחרון. יומן התצפיות הוא המערך
\[
 H_{S,\tau}(b)=\bigl(b_k(i)\bigr)_{0\le k\le\tau,\,i\in S}.
\]
בקובץ \textenglish{\nolinkurl{observer/Rule110Observer.lean}}, הפרדיקט \textenglish{\texttt{CanPredict}} דורש מפענח שמפיק את $b_t(s)$ מתוך היומן הזה בכל תצורה קבילה. המשפט \textenglish{\texttt{canPredict\_iff}} קובע כי דרישה זו שקולה לכך ש־$b_t(s)$ קבוע על הסיבים של $H_{S,\tau}$. ההגדרה מאפשרת להטיל מגבלות כלשהן על התצורות ההתחלתיות.

נסמן ב־$I_t(s)=s+I_t$ את תחום התלות המדויק לאחר הזזה. נשתמש במשפט \textenglish{\texttt{Rails.essential\_iff}}, הקובע כי $D_t=I_t$. אינווריאנטיות הדינמיקה תחת הזזה מעבירה את התוצאה מראשית הצירים אל $s$. הפרדיקט \textenglish{\texttt{CoversTarget}} פירושו
\begin{equation}
 I_t(s)\ \subseteq\ \bigcup_{i\in S}\bigcup_{k=0}^{\tau} I_k(i).
 \label{eq:observer-cover}
\end{equation}
עבור מחלקת כל התצורות ההתחלתיות זהו תנאי הכרחי. אם אתר $j$ שבאגף שמאל נעדר מאגף ימין, מהותיותו של $j$ מבטיחה את קיומן של שתי תצורות הנבדלות רק ב־$j$ ומפיקות ערכים שונים של $b_t(s)$. הקריאות ביומנים שלהן זהות, מפני ש־$j$ נמצא מחוץ לתחומי התלות של אותן קריאות. זוג זה מפר את קריטריון החיזוי. הבחירה במחלקת המצבים ההתחלתיים חשובה: במחלקה בת תצורה אחת אפשר לחזות כל סיבית עתידית באמצעות ערך קבוע.

\paragraph{דוגמה נגדית למספיקות התנאי.}
נבחר $S=(-1,1)$, $\tau=1$, $s=0$ ו־$t=2$. חרוטי התלות של הקריאות האחרונות לבדן הם $[-2,0]$ ו־$[0,2]$, והם מכסים את $I_2(0)=[-2,2]$. בתצורה ההתחלתית הראשונה האתר היחיד שערכו אחת הוא $-1$. בתצורה השנייה הערך אחת מופיע באתרים $-1$ ו־$0$. בשתי התצורות ערכי כל שאר האתרים הם אפס. אחרי צעד אחד קבוצות האתרים שערכם אחת הן $\{-2,-1\}$ ו־$\{-2,-1,0\}$, בהתאמה. בשני העולמות החיישנים מפיקים אותו יומן $(10,10)$. בעדכון השני באתר $0$, השלשות המקומיות $100$ ו־$110$ נותנות $0$ ו־$1$, בהתאמה. לכן תנאי הכיסוי מתקיים, אך החיזוי מתוך היומן המלא אינו אפשרי. הבניות המדויקות מופיעות ב־\textenglish{\texttt{separated\_cells\_counterexample}} וב־\textenglish{\texttt{coverage\_does\_not\_suffice}}.

התופעה מתרחשת גם בחיישנים סמוכים. עבור $S=(0,1)$ ו־$\tau=2$, החרוטים $[-2,2]$ ו־$[-1,3]$ מכסים את $I_3(0)=[-2,3]$. קבוצות האתרים ההתחלתיות שערכם אחת, $\{2\}$ ו־$\{-1,2\}$, נותנות אותו יומן $(00,01,11)$ ואת הערכים המנוגדים $1$ ו־$0$ באתר $0$ בזמן $3$. זו הדוגמה \textenglish{\texttt{adjacent\_cells\_counterexample}}. יש גם תנאי מספיק שמשמעותו ישירה: אם השורה הנצפית האחרונה נותנת את ערכי \emph{כל התאים} ב־$I_{t-\tau}(s)$ כאשר $t\ge\tau$, ההתפתחות הרגילה מאותם ערכים קובעת את היעד באופן יחיד. תנאי~\eqref{eq:observer-cover} אינו שולל אובדן יכולת הבחנה כאשר הנתונים ההתחלתיים נדחסים לקומץ קריאות של תאים.

\subsection{ארבע היסטוריות מוצלחות ובקרת 11 סיביות ממצה}

נניח מילה בת 11 סיביות על הישר השלם באתרים $0,\ldots,10$, משמאל לימין, על רקע אפסים. לפני העדכון הראשון נכתוב אחת באתר $11$. סיבית היעד מן העבר היא $Y=b_0(5)$, קריאת החיישן היא $Z=b_2(10)$, יומן המקבל הוא $H=(b_k(0))_{k=0}^3$, והסיבית האחרונה בו היא $K=b_3(0)$. אחרי התצפית בזמן 3 נכתוב את $Z$ באתר $-1$ ונקרא את תוצאת העדכון הרגיל של תא זה בזמן 4.

\begin{center}
\begin{tabular}{lccccc}
מילה התחלתית & $H$ & $K$ & $Y$ & $Z$ & פלט בזמן 4 \\
\hline
\textenglish{\texttt{10111010001}} & 1100 & 0 & 0 & 0 & 0 \\
\textenglish{\texttt{10111110111}} & 1100 & 0 & 1 & 1 & 1 \\
\textenglish{\texttt{11000010111}} & 1101 & 1 & 0 & 1 & 0 \\
\textenglish{\texttt{11001100001}} & 1101 & 1 & 1 & 0 & 1
\end{tabular}
\end{center}

בארבע ההיסטוריות האלה מתקיים $Y=K\oplus Z$. בכל זוג שבו $K$ קבוע, אפילו התצורה המלאה של שכבת היסוד בזמן 3 זהה לפני הכתיבה. ערך שכנו השמאלי של התא המקבל הוא אחת, וערך שכנו הימני הוא $K$. לכן הסיבית המוזרקת נותנת את התשובה הנכונה על פי הזהות
\begin{equation}
 f_{110}(1,z,k)=z\oplus k.
 \label{eq:observer-xor-gate}
\end{equation}
היפוך הסיבית שנמסרה הופך את התשובה בכל ארבעת המקרים. אלא ששתי המילים שהוכנו בכל זוג נבדלות גם באתרים אחרים מלבד $5$.

בכל המחלקה האחידה של $2^{11}$ המילים נשווה כל מילה לעותק שלה שהתקבל מהיפוך $b_0(5)$ בלבד. החיישן באתר $10$ בזמן 2 תלוי באתרים ההתחלתיים $8,\ldots,12$. יומן התא $0$ עד זמן 3 נקבע כולו על ידי אתר $0$ וסביבת שלושת הצעדים שלו. לפיכך, בכל אחד מ־1024 הזוגות, שני ערכי $Y$ נותנים אותן תצפיות $(H,Z)$. כל תשובה בינרית דטרמיניסטית המבוססת על $(H,Z)$ נכונה בדיוק ב־1024/2048 מן ההיסטוריות. מנייה בקובץ \textenglish{\nolinkurl{observer/domain.py}} מאשרת 0/1024 שינויים בקריאת החיישן ואת החסם 1024/2048 למפענח. הטענה על ההשפעה הסיבתית של הסיבית \emph{המוזרקת} על התא המקבל נפרדת מן הטענה על השפעת סיבית \emph{העבר} על החיישן.

\subsection{היפוך יחיד ושחזור בזמן 5}
\label{sec:observer-return}

כעת נניח כל מילה שבטבלה הבאה באתרים $-6$ עד $13$, כולל הקצוות, על רקע אפסים מחוץ לקטע. בבנייה זו הצופה קורא באופן פסיבי את $Z=b_2(3)$. לפני ההחזרה מתקיים $Y=b_0(5)$ וכן $K=b_4(0)$. לכל ערך קבוע של $K$, שתי המילים נבדלות רק באתר $5$.

\begin{center}
\begin{tabular}{cc l cc}
$K$ & $Y$ & המילה על $[-6,13]$ & $Z$ & פלט בזמן 5 \\
\hline
0 & 0 & \textenglish{\texttt{00000001100001000000}} & 0 & 0 \\
0 & 1 & \textenglish{\texttt{00000001100101000000}} & 1 & 1 \\
1 & 0 & \textenglish{\texttt{00010111010001000010}} & 1 & 0 \\
1 & 1 & \textenglish{\texttt{00010111010101000010}} & 0 & 1
\end{tabular}
\end{center}

היפוך של $Y$ בלבד משנה את $Z$ עבור שני ערכי $K$. כאשר $K=0$, החיישן מחזיר את $Y$; כאשר $K=1$, הוא מחזיר את הערך המשלים שלו. בזמן 4 תצורות שכבת היסוד בכל זוג זהות לאורך הישר כולו. גם היומנים שלהן בחלון הקליטה $(-2,-1,0)$ זהים בכל הזמנים מ־0 עד 4 לפני הכתיבה. בנקודה זו $b_4(-2)=1$. כתיבה חד־פעמית של $Z$ באתר $-1$ ועדכון אחד לפי כלל~110 נותנים $b_5(-1)=Z\oplus K=Y$, על פי~\eqref{eq:observer-xor-gate}. כתיבת $1-Z$ במקום זאת נותנת $1-Y$ בכל ארבעת העולמות. לפני הכתיבה שכבת היסוד כבר איבדה את ההבחנה, ואילו האוגר החיצוני בן הסיבית האחת שימר אותה. החישוב על הישר האינסופי ממומש בקובץ \textenglish{\nolinkurl{observer/causal_line.py}}.

\paragraph{תחולת תוצאת המינימליות.}
זמן 4 הוא המועד המוקדם ביותר להתלכדות מלאה במחלקה הבאה על הישר. בכל זוג הופכים תא התחלתי אחד בלבד. נלקחת קריאה גולמית אחת באתר קבוע ובזמן שנקבע מראש. שני הקלטים הבינריים $K,Z$ חייבים להיות נחוצים יחד לשחזור מדויק על פני הריבוע של ארבע ההיסטוריות. כדי לבדוק התלכדות עד זמן 3, די להביא בחשבון את כל ההבדלים האפשריים באותו זמן; כולם נמצאים ברדיוס 3 מן האתר שהפכנו, ונקבעים על ידי קטע התחלתי ברדיוס 6. בלי התא שהפכנו נותרות 12 סיביות סביבו. מתוך $2^{12}=4096$ הסביבות, שמונה מובילות להתלכדות מלאה עד הצעד השלישי. נקבע חיישן רגיש כלשהו וזמן קריאה לפני ההתלכדות. בכל אותן סביבות שבהן הוא רגיש, כיוון תגובתו להיפוך זהה. לשם שחזור מדויק החיישן צריך להיות רגיש בשני ההקשרים. אילו כיוון התגובה היה זהה, $Y$ היה פונקציה של $Z$ בלבד, בניגוד לדרישה ששני הקלטים יהיו נחוצים יחד. ארבע המילים שבטבלה משיגות את חסם ההתלכדות ומפיקות את הפלט בעדכון הבא. הסביבות נמנות בקובץ \textenglish{\nolinkurl{observer/causal_crosscheck.py}}.

\subsection{מסירה במהירות סופית וביצוע אוטונומי}

בניסוי הקודם התקן חיצוני כתב את הסיבית שנשמרה ישירות באתר הקליטה. עם אותן ארבע מילים התחלתיות אפשר לחסום במפורש את מהירות הערוץ. בזמנים $2,3,4,5,6,7$ הסיבית המבודדת נמצאת באתרים $3,3,3,2,1,0$, בהתאמה. לוח הזמנים זהה בכל העולמות, והסיבית נעה לכל היותר תא אחד בצעד. בזמן 7 ערכי השכנים של תא $0$ הם $b_7(-1)=1$ ו־$b_7(1)=K$. לאחר כתיבת הסיבית שנמסרה $Z$, עדכון רגיל נותן $b_8(0)=Y$. החלפת $Z$ בערך המשלים שלו הופכת את כל ארבעת הפלטים. לפני המסירה, יומני החלון $(-1,0,1)$ זהים בכל זוג שבו $K$ קבוע. הבדיקה ממומשת בקובץ \textenglish{\nolinkurl{observer/local_return.py}}.

המימוש האוטונומי משתמש בטבעת באורך 61 ובמעבר מקומי הומוגני אחד שאינו תלוי בזמן, ברדיוס 1. כל תא שומר סיבית יסוד, מונה מופע, תפקיד קבוע וסיבית מסר אחת. באתחול מקצים את תפקיד הצופה לאתר 3 ואת תפקיד המקבל לאתר 0. המופע שולט בלכידת המסר, בהעברתו המקומית ובהכנסתו. לאחר האתחול אין פקודות חיצוניות. בכל זוג עולמות, שכבות היסוד זהות בזמנים 4 עד 7, ואילו המצבים המורחבים המלאים נשארים שונים. הבדיקה בקובץ \textenglish{\nolinkurl{observer/autonomous.py}} נותנת 4/4 פלטים נכונים בזמן 8. קיומו של אלפבית סופי מורחב כזה משאיר פתוחה את \textenglish{R110.11}, המבקשת את אותו פרוטוקול מלא במסגרת כלל~110 יחיד, הומוגני ובינרי ממש. לאחר שהמצבים המלאים של שתי מערכות דטרמיניסטיות מבודדות בעלות אותה דינמיקה נעשים זהים, הבחנה יכולה להופיע מחדש רק אם מתקבל קלט חדש המבדיל ביניהן.

\paragraph{מצב האימות הפורמלי.}
בגרסה ההיסטורית, גרסה נפרדת של מודול הצופה ב־\textenglish{Lean~4.33.1} הפיקה תשעה דוחות \textenglish{\texttt{\#print axioms}}, ובהם קריטריון הסיבים לחיזוי ודוגמאות נגדיות למספיקות הכיסוי. במסגרת בניית המונוגרפיה, \textenglish{Lean~4.35.0} בדק את המקור הקנוני עם הייבואים \textenglish{\texttt{Rule110Exact}} ו־\textenglish{\texttt{Rule110FrontGlue}} והחזיר את כל 11 רשימות האקסיומות שהתבקשו. נספח~\ref{app:lean} מציין את יומני הבדיקה. המסלולים ותוצאות המנייה מספקים עדויות נפרדות לדוגמאות הסופיות בסעיף זה.

\section{העברת כלל בסיבית אחת ב־\textenglish{OLMo} שמשקליו מוקפאים}
\label{sec:olmo}
בניסוי בכלל 110, לאחר שהתצורות מתלכדות, הצופה כותב בחזרה לתא את קריאת החיישן המקורית $Z$. בעדכון הבא האוטומט משחזר את הסיבית הקודמת $Y$ בעזרת הסיבית $K$ שנותרה בסריג. בניסוי הדו־שלבי ב־\textenglish{OLMo}, לפני שמוצגת מחרוזת חדשה, השולח מצמצם את בחירתו הקודמת $h$ ואת המשוב $e$ לטוקן יחיד המייצג כלל, $m\in\{A,B\}$. המקבל משתמש באותם משקלים מוקפאים, פותח אפיזודה חדשה עם מטמון \textenglish{KV} חדש ומחיל את הכלל על המחרוזת. הגבול בין האפיזודות מאפשר למדוד העברה של כלל שכבר נבחר. כדי למדוד תיקון פנימי של אסטרטגיה תוך שמירה על רציפות מצב המודל נדרש ניסוי נפרד.

\subsection{הריבוע הנגד־עובדתי}
נניח שהמודל בחר קודם לכן $h=0$ (\textenglish{A}, לקחת את הסיבית הראשונה) או $h=1$ (\textenglish{B}, לקחת את הסיבית האחרונה). הצופה מחזיר $e=0$ אם הבחירה נכונה ו־$e=1$ אם שגויה. שני ענפי המשוב זמינים לכל בחירה. לפני הצגת מחרוזת חדשה, הכלל המתוקן הוא $p=h\oplus e$. לאחר הצגת $\mathbf{x}=LwR$, התשובה הנכונה היא $f_0(\mathbf{x})=L$ או $f_1(\mathbf{x})=R$, בהתאם לכלל. כאשר סיביות הקצה שונות זו מזו, שתי ההיסטוריות $h$ עם אותם $e,\mathbf{x}$ מחייבות פעולות הפוכות. כאשר סיביות הקצה זהות, שתיהן מחייבות אותה סיבית. הצגת כל ארבעת הצירופים של סיביות הקצה משמשת בקרה נגד אסטרטגיה שלמדה רק לקשר בין התשובה לבין תווית המשוב.

כל מחרוזת אמצע מגדירה בלוק של 16 תנאים $(L,R,h,e)$. מערך הפיתוח כולל שמונה מחרוזות אמצע שנבחרו לפני המדידה. בלוק נחשב הצלחה רק אם כל 16 מתוך 16 התשובות נכונות. סף המעבר הוא 6 מתוך 8 בלוקים. בכל קריאה מספרית נדרשת הסתברות כוללת של $P(0)+P(1)\ge0.80$ ביחס לכל אוצר המילים. סימן ההפרש בין הלוגיטים $\ell_1-\ell_0$ קובע איזו סיבית נבחרה כאשר ההפרש אינו אפס. הפרש אפס נחשב להימנעות מהכרעה. בשורות בעלות אותה תחילית נעשה שימוש חוזר בחישוב המודל המתאים.

\subsection{מדידות}
ב־\textenglish{G0} מבודדים את המטמון לאחר עיבוד הטוקן שנבחר ולפני כל קריאה אבחונית. ב־\textenglish{G0R} מחליפים את הערכים בעמדת הבחירה בערכים מההיסטוריה האחרת, שהתקבלה באופן טבעי, ואז משחזרים את הערכים המקוריים. בהרצה המקורית של \textenglish{Base}, כל 256 מתוך 256 ההשוואות של וקטורי הלוגיטים המלאים בשני כיווני ההחלפה שחזרו את הבתים המקוריים, וההפרש המוחלט המרבי היה אפס. בקרה זו בודקת את תקינות השחזור. השאלה אם המודל בוחר באסטרטגיה ההפוכה היא שאלה התנהגותית נפרדת.

\begin{center}\small
\begin{tabularx}{\linewidth}{@{}Xrrrl@{}}\toprule
גרסת משקלים וממשק (פיתוח) & שורות & בלוקים & זוגות כאשר $e=1$ & מצב \\\midrule
\textenglish{Base}, ממשק גולמי מקורי & 76/128 & 0/8 & 0/16 & \textenglish{\texttt{FAIL\_DEV}}\\
\textenglish{Instruct}, ממשק גולמי מותאם & 79/128 & 0/8 & 1/16 & \textenglish{\texttt{FAIL\_FORMAT}}\\
\textenglish{Instruct}, שיחה, \textenglish{correct/wrong} & 115/128 & 2/8 & 9/16 & \textenglish{\texttt{FAIL\_DEV}}\\
\textenglish{Instruct}, שיחה, \textenglish{keep/switch} & 109/128 & 0/8 & 6/16 & \textenglish{\texttt{FAIL\_DEV}}\\
\textenglish{Instruct}, שיחה, שני שלבים (\textenglish{A}) & 124/128 & 6/8 & 16/16 & \textenglish{\texttt{PASS\_DEV}}\\
\bottomrule\end{tabularx}
\end{center}
כדי שזוג ייחשב הצלחה בטבלה נדרשות שתי תשובות נכונות להיסטוריות מנוגדות כאשר $e,\mathbf{x}$ זהים ומתקיים $L\ne R$. בין שתי השורות הראשונות הוחלפו הן המשקלים הן הדוגמאות שהוצגו למודל. בממשק הגולמי של \textenglish{Instruct}, ב־28 מתוך 64 שורות הבקרה שבהן ניתן כלל מפורש, מסת ההסתברות הייתה מתחת לסף. בתבנית השיחה המקורית של המודל, עם המשוב \textenglish{correct/wrong}, קריאות אבחוניות בעותקים נפרדים של המטמון זיהו נכונה 2 מתוך 2 בחירות קודמות. כשניתן כלל מפורש, המודל פתר 62 מתוך 64 שורות ו־6 מתוך 8 בלוקים. מקר ועמיתיו מדדו בנפרד אם מודל יכול לזהות התערבות באקטיבציות שלו ואם הוא יכול לנקוב בשם המושג שהוכנס באמצעות הווקטור המוזרק~\cite{macar2026}. הקריטריון ההתנהגותי שלנו מחייב את המודל להחיל את המידע שקיבל על מחרוזת חדשה. ניסוי \textenglish{G3} המבוסס על התערבות באקטיבציות נושא את המצב \textenglish{\texttt{NOT\_RUN}}.

\subsection{משפט על ערוץ של סיבית אחת}
בגרסה הדו־שלבית השולח מפיק, לפני הצגת $\mathbf{x}$, את הטוקן היחיד שציונו מרבי בכל אוצר המילים. רק הטוקנים \textenglish{A} ו־\textenglish{B} קבילים. המקבל מתחיל עם מטמון \textenglish{KV} ריק משלו. הוא מקבל רק את מזהה הטוקן $m$, מחשב מחדש את ייצוגו, ואז מקבל את $\mathbf{x}$. מזהה הענף של השולח, האקטיבציות והלוגיטים שלו אינם נגישים למקבל. התוכנית החיצונית אינה משתמשת בתוויות הכללים כדי לבחור את הקוד.

\begin{proposition}
נניח כי $m$ הוא הקלט היחיד למקבל שיכול להשתנות בהתאם לאפיזודה הקודמת, כי לאלפבית שלו שני איברים, וכי הערוץ הוא מהצורה $(h,e)\longrightarrow m\longrightarrow F(m,\mathbf{x})$. אם עבור קלט יחיד המשותף לכל ההיסטוריות ומקיים $L\ne R$ מתקבל פתרון מדויק בכל ארבע ההיסטוריות, אזי טבלת השולח, לפי הסדר $00,01,10,11$, היא $0110$ או $1001$.
\end{proposition}
\begin{proof}
עבור $\mathbf{x}$ הנתון, שני הערכים של $p$ מחייבים תשובות שונות. שיוך אותו $m$ להיסטוריות בעלות $p$ שונה היה מציג למקבל קלטים זהים עבור יעדים שונים. לכן שתי קבוצות המסרים, שאינן ריקות, מוכרחות להיות זרות זו לזו. יש רק שני מסרים, ולכן לכל $p$ מותאם בדיוק $m$ אחד. אפשריות אפוא רק שתי התמורות של תוויות הטוקנים.
\end{proof}

הטבלה שנצפתה בממצאי הפיתוח הייתה $A,B,B,A$. לכל $e$, שתי ההיסטוריות נפתרו בכל 16 מתוך 16 הזוגות שבהם סיביות הקצה שונות. התוצאה הייתה 64 מתוך 64 תשובות נכונות בעת מסירת הטוקן, ו־32 מתוך 64 בבקרת \textenglish{A} קבוע שחושבה בנפרד. ארבע החלטות השולח משמשות מחדש במחרוזות החדשות של המקבל. לפיכך המסר מכיל מידע על מחלקת הכלל עוד לפני שהמחרוזת מוצגת. כדי לברר כיצד השולח מחשב את המידע הזה נדרשות התערבויות אחרות. מצב רב־ממדי יכול להעביר את שתי הסיביות $(h,e)$ ולהותיר למקבל את חישוב הכלל. נספח~\ref{app:gates} מפרט את מקורות המדידות ואת הבקרות.

\section{אי־הכרעה בדבר חוויה וגשר מותנה}
\label{sec:boundary}
שמירת סיבית מבססת תכונה של מערכת ביחס להתערבויות מוגדרות. האם אפשר להסיק מן התצפיות האלה שלמערכת מסוימת זו, ברגע המסוים הזה, יש חוויה סובייקטיבית?

\paragraph{\textenglish{T5}.} נניח שכל האקסיומות של התאוריה האופרציונלית $T$ מנוסחות ללא פרדיקט החוויה $E(a,t)$. ניקח מודל כלשהו $M$ של התאוריה, ובו נקודה מסומנת $u=(a,t)$. נרחיב אותו תחילה באמצעות פירוש $E$ כקבוצה הריקה, ואחר כך באמצעות פירוש $E$ כקבוצה המכילה את $u$. מכיוון שרק הפרדיקט החדש משתנה, שני המבנים מקיימים את $T$ ומסכימים בתוצאותיהן של כל ההתערבויות האופרציונליות. במבנה הראשון $E(u)$ שקרי, ובשני הוא אמיתי. לכן $T$ אינה קובעת אף אחד משני ערכי האמת, גם כאשר הקריטריון התפקודי לעדכון מתקיים. זוהי תוצאה על כוח הביטוי של התאוריה. בניית שני המבנים מותירה פתוחה את השאלה אם שתי חוויות שונות ניתנות למימוש פיזי באותו מצע.

\paragraph{\textenglish{T6}.} לחבילה $B=(B_{\mathrm{id}},B_{\mathrm{nom}},B_{\mathrm{epist}})$ ניתן המצב \textenglish{\texttt{DECLARED\_NOT\_DERIVED}}. שלושת חלקיה מגדירים את הישות שעליה מוסב ההיסק, את ההנחה הנוספת ואת הכללים לפירוש הראיות.
\begin{description}[leftmargin=1.35cm,style=nextline]
\item[$B_{\mathrm{id}}$] נגדיר את גבול המערכת, מאקרו־מצב וחבורה של שינויי תיוג קבילים של הקואורדינטות. התיאור אינו משתנה תחת שינויי תיוג אלה. בהינתן קלט חיצוני קבוע, על המצב להשפיע על המצב העוקב לו דרך שרשרת סיבתית פנימית. שושלת רציפה מחייבת מצב עוקב יחיד. שכפול סימטרי של מטמון \textenglish{KV}, שלאחריו שני העותקים ממשיכים באופן בלתי תלוי, נרשם כאירוע \textenglish{\texttt{BranchEvent}}. לשני הצאצאים הסיבתיים אב קדמון משותף ומזהי ענף נפרדים. זוהי מוסכמה למעקב אחר היסטוריות. שאלת הזהות המספרית של הסובייקט החווה נותרת פתוחה.
\item[$B_{\mathrm{nom}}$] נכתוב $C=I_{\mathrm{cont}}\land C_{\mathrm{med}}\land A_{\mathrm{ind}}\land D_{\mathrm{dyn}}$. כאן $I_{\mathrm{cont}}$ מגדיר רציפות. $C_{\mathrm{med}}$ מחייב תיווך פנימי כאשר הגישה לעמדות שבהן נמצאים הבחירה הקודמת והאות חסומה, וכן העברה של אותו כלל לסוג משימה אחר. $A_{\mathrm{ind}}$ מחייב בדיקות בלתי תלויות שלפיהן כמה מודולים יכולים לגשת לתוכן משותף, והערכת מהימנות יכולה לשנות פעולה. $D_{\mathrm{dyn}}$ מחייב התארגנות מחדש מתמשכת, הנגרמת על ידי העיבוד ומשפיעה סיבתית על המשך העיבוד לאחר האפיזודה. נאמץ כהשערה את התנאי המספיק $\forall u\,[C(u)\Rightarrow E(u)]$.
\item[$B_{\mathrm{epist}}$] המצב \textenglish{\texttt{UN}} מציין שאלה שנותרה פתוחה לנוכח הראיות הקיימות. כל אחת מתוויות ההחלטה \textenglish{\texttt{NONE}}, \textenglish{\texttt{INVESTIGATION\_PRIORITY}} ו־\textenglish{\texttt{SENTIENCE\_CANDIDATE}} מחייבת נימוק נפרד ומפורש. לכימות דרגת האמונה דרושים הסתברויות אפריוריות ומודל תצפית.
\end{description}
\begin{theorem}[משפט מותנה \textenglish{T6}]
\label{thm:conditional-bridge}
בהינתן $B_{\mathrm{nom}}$, מן $C(u)$ נובע $E(u)$. אם קיים לפחות מודל אופרציונלי אחד $M$ העולה בקנה אחד עם הראיות ועם $B_{\mathrm{id}}$, שבו $C(u)$ שקרי, אז אותן ראיות מתיישבות עם שני ערכי האמת של $E(u)$ גם לאחר הוספת $B_{\mathrm{nom}}$.
\end{theorem}
\begin{proof}
לטענה הראשונה נחיל את הגרירה בנקודה $u$. לטענה השנייה נגדיר $S=\{v:C(v)\}$ ב־$M$, ונרחיב את אותו מודל אופרציונלי באמצעות הפירושים $E_0=S$ ו־$E_1=S\cup\{u\}$. שניהם מכילים את $S$ ומקיימים את הנחת הגשר. מכיוון ש־$u\notin S$, הם נבדלים בערך של $E(u)$. שתי ההרחבות חולקות אותם נתונים אופרציונליים ואותן תוצאות בכל אחת מההתערבויות המוגדרות.
\end{proof}
אי־מדידה של אחד מרכיבי הצירוף אינה מוכיחה כשלעצמה שאותו רכיב שקרי. במקרה של כלל 110, אוגר בן סיבית אחת, מפענח המבוסס על טבלת חיפוש ו־\textenglish{OLMo} עם משקלים מוקפאים, המצב האופרציונלי של $E(a,t)$ בעבודה זו הוא \textenglish{\texttt{UN}}. העברה באמצעות טוקן חיצוני אינה מכריעה בשאלת התיווך הפנימי ובשאלת הסמנים הארכיטקטוניים הבלתי תלויים. מסגרת ההפשטה הסיבתית של גייגר ועמיתיו~\cite{geiger2025} ישימה לבדיקה המוצעת \textenglish{AI.04}. חיפוש התאמה מבוזרת (\textenglish{DAS}) מזהה תת־מרחב מועמד עבור המשתנה $p$. בדוגמאות בלתי תלויות, כאשר הגישה לעמדות המקוריות $h,e$ חסומה, אפשר להציב בתת־המרחב את הערך מהיסטוריית תורם ולבדוק אם התשובה נקבעת לפי הכלל של התורם. גם שינוי במטמון \textenglish{KV} במהלך אפיזודה אינו מספיק כדי לבסס פלסטיות מתמשכת.

\section{דיון וכיוונים לבדיקות נוספות}
בשני הניסויים, בקרת אבחון מועילה משתמשת בשתי היסטוריות שבהן אותה תצפית מחייבת תשובות שונות. מערכת שרואה רק את המשותף להיסטוריות אינה יכולה להשיב נכונה בשתיהן. זוג כזה קובע איזה מידע נחוץ. כדי לייחס למערכת מסוימת שימוש במידע הזה, יש להראות גם מה קורה כאשר הנשא שלו נמסר, מוחלף או מוסר, תוך שמירה על קלט ההמשך ללא שינוי. כיסוי גאומטרי של חרוט תלות, משתנה מתואם שניתן לשחזר מן הנתונים ודיוק ממוצע גבוה אינם מעידים כשלעצמם על השפעה כזאת.

\paragraph{תחולת המסקנות.}
הקשר בין קידוד ללא שגיאות לצביעת גרפים מוכר בתורת המידע \cite{witsenhausen,charpenay}. התרומה כאן היא בהגדרת התערבויות הניתנות לבדיקה, במתן עדויות בדידות מפורשות ובהגבלת הערוץ בין שתי אפיזודות. על הישר האינסופי, לכל זמן $t$ ולכל קואורדינטה $j$ בתוך החלון קיים זוג תצורות הנבדלות רק ב־$j$ ונותנות סיביות יעד שונות. מחוץ לחלון, שינוי $j$ בלבד משאיר את סיבית היעד ללא שינוי בכל תצורה התחלתית. בניסויי \textenglish{OLMo}, כל שורה בטבלה נוגעת לשילוב מסוים של גרסת משקולות קפואה (\textenglish{Base} או \textenglish{Instruct}) וממשק במערך בדיקה סופי אחד. השאלה אם למערכת יש חוויה סובייקטיבית או אין חוויה כזאת חורגת ממה שאפשר להסיק מן הנתונים האלה.

\paragraph{יישום בתחומים אחרים.}
בבדיקה מנגנונית יש לוודא תחילה שהמודל מסוגל לפתור את המשימה כשההיסטוריה נשמרת במלואה. לאחר מכן אפשר להעביר נשא חבוי מועמד ולבדוק מה קורה כשמוחקים אותו. בחקר התודעה, ובכלל זה בדיון על הפרעות בתודעה, מועיל להבחין בין זמינות של דיווח, הנחות בדבר חוויה והחלטה לנקוט אמצעי זהירות. יישום קליני של ההבחנה הזאת דורש אימות בלתי תלוי. שאלת מקורו של מצב עולה גם בחקירה פורנזית דיגיטלית ובמחשוב מבוזר. הקלטה נכונה יכולה לשחזר תשובה חיה לשאילתה קבועה. בסדרת בדיקות נפרדת שבה הכתובת השתנתה, הערוץ החי השיג 8/8 וההקלטה המוקדמת השיגה 6/8. אפשר לבדוק בנפרד את מקורה של הסתעפות ואת יכולתה להשיב לשאילתה חדשה.

\paragraph{משוב ממוען ב־\textenglish{OLMo (T3, \texttt{NOT\_RUN})}.}
עם סיום הבדיקה במערך 33--44, ניסוי זה ידרוש מערך בדיקה נפרד. המודל שומר שתי בחירות קודמות שלו $h_0,h_1$. עליו להסיק מתוכנה של משימה חדשה איזו בחירה טעונה תיקון. כתובת הייחוס $K$ מוסתרת. לפני שהמודל רואה את המחרוזת $\mathbf{x}$, הוא מפיק בעצמו סיבית שאילתה $q$. הצופה מחזיק שתי סיביות משוב $e_0,e_1$ ומחזיר אחת מהן. לאחר מכן על המודל להחיל את כלל היעד $h_K\oplus e_K$ על המחרוזת. תשובת הצופה היא מקור המידע היחיד של המודל על סיביות המשוב.

עבור אותן היסטוריות, מוסרים לצופה את הכתובת $q$ (שאילתה חיה), $0$ (כתובת קבועה) או $1-q$ (כתובת מוחלפת). בכל תנאי המקבל רואה סיבית אחת באותו מיקום. מתעדים את הכתובת שנמסרה ואת החיווי אם בוצעה החלפה, ומסתירים את שניהם מהמקבל. מודדים את הבחירה בשאילתה המועילה ($q=K$) בנפרד מדיוק התשובה הסופית כאשר $L\ne R$. השוואה מזווגת כאשר $e_0\ne e_1$ מראה אם התשובה השתנתה בעקבות החלפת הכתובת.

במחרוזות שקצותיהן שונים, שני ערכי $K$ מופיעים באותה תדירות. לכל משימה, זוג בחירות עבר ומחרוזת, כל ארבעת זוגות המשוב מופיעים באותה תדירות. תקרת הדיוק בכתובת קבועה היא $75\%$. כאשר $K=0$, התשובה $e_0$ נותנת את המשוב הדרוש. כאשר $K=1$, שני ערכי $e_1$ נשארים שווי הסתברות גם לאחר קבלת $e_0$. לכן הדיוק המיטבי הוא $\tfrac12\cdot1+\tfrac12\cdot\tfrac12=\tfrac34$. השאלה אם שאילתה חיה חורגת מן החסם הזה נבחנת במערך המאוזן המלא של המחרוזות שקצותיהן שונים.

\begin{center}\small
\begin{tabularx}{\linewidth}{@{}p{.23\linewidth}X@{}}\toprule
משימה בגרסה הנוכחית & ראיות נדרשות \\\midrule
\textenglish{G4-1}, מערך 33--44 & יש לקבע את שני הממשקים המושווים על מערך הפיתוח לפני הפתיחה הראשונה של מערך הבדיקה המשותף. העברה מדויקת בסיבית אחת מחייבת השלמה מוצלחת של כל 12/12 הבלוקים, וכן את בקרות המחיקה, הפורמט וגבולות הערוץ.\\
\textenglish{G3}, מערך 21--32 & אחרי שהערוץ הטקסטואלי עובר את מבחני הסף, יש להחליף את הטקסט בפעימת אקטיבציה אחת שנקבעה מראש. כיבוי הפעימה צריך להפוך את שני ענפי המשוב לזהים עבור אותה בחירה ואותו קלט.\\
\textenglish{AI.04}: \textenglish{G4} רב־ממדי ו־\textenglish{DAS} & יש להעביר מצב פנימי מוגדר כשאין גישה למיקומי המקור. כדי לאתר תת־מרחב, יש להגדיר מראש את השכבה, מדד התיווך, סף הסבילות ובקרה עצמאית. גם $C_{\rm med}$ מחייב משימה נוספת המשתמשת באותו כלל.\\
\textenglish{R110.11}, משבצת \textenglish{\texttt{leaf:A6}} & יש לממש את האחסון וההחזרה במלואם בכלל~110 בינרי והומוגני. הבנייה האוטונומית שהודגמה משתמשת באלפבית מורחב.\\
פורמליזציה של ההתלכדות וההחזרה & יש לבצע פורמליזציה של הבדיקה הממצה של 4096 הסביבות ושל העדות הקונסטרוקטיבית עבור $t=4\to5$ ב־\textenglish{Lean}. האימות של משפט הגאומטריה היסודי ושל 11/11 השאילתות במודול הצופה שוחזר (נספח~\ref{app:lean}).\\
$D_{\rm dyn}$ & יש להראות שינוי בארגון שנמשך מעבר לאפיזודה, ואת השפעתו על קלט חדש בתנאי איפוס מבוקרים.\\
\bottomrule\end{tabularx}
\end{center}
פרוטוקול השלבים בנספח~\ref{app:gates} מפרט באילו תוצאות הערוץ המוצע נכשל במבחן. תוצאה שלילית עבור ערוץ מסוים מגבילה את מה שאפשר להסיק על הנשא שנבחר. נספח~\ref{app:registers} מונה את ההשערות הקשורות ואת מצבן.

\appendix

\section{אימות פורמלי של הגאומטריה ויישומה לבעיית הצופה}\label{app:lean}

\subsection{שני כיווני ההוכחה}
הקובץ \textenglish{\texttt{lean/Rule110Exact.lean}} מייבא את \textenglish{\texttt{Finite}} ואת \textenglish{\texttt{Late}}. משפט השורש נקרא \textenglish{\texttt{essential\_iff}} במרחב השמות \textenglish{\texttt{Rails}}. אגפו השמאלי, \textenglish{\texttt{EssentialAt}}\,$t\,j$, דורש שתי תצורות הזהות בכל מקום למעט, אולי, בקואורדינטה $j$, והמניבות ערכים שונים בתא המטרה בזמן $t$. ההבדל בפלט מחייב שהקלטים יהיו שונים ב־$j$. באגף הימני, \textenglish{Lean} מחשבת תחילה את $t-t/3$ במספרים הטבעיים באמצעות חילוק שלם, ואז משכנת את התוצאה ב־$\mathbb Z$.

בכיוון הראשון, \textenglish{\texttt{left\_dependency\_bound}} ו־\textenglish{\texttt{right\_dependency\_bound}} שוללים קואורדינטות מהותיות מחוץ לחלון. בכיוון ההפוך, התעודה מספקת זוג מסלולים קביל לכל $(t,j)$ בחלון. עבור $t<2485$, הקובץ \textenglish{\texttt{lean/Finite.lean}} מחיל את \textenglish{\texttt{checkSound}} על המקטעים הסופיים, מחלץ סיבית המשמשת עדות במיקום הפיזי של הפגם $s=-j$, ומעביר את העדות לישר האינסופי. המחולל \textenglish{\texttt{scripts/generate.py}} מספק תמונות מצב מועמדות בצורת וקטורי סיביות. \textenglish{Lean} בודקת התאמה בגבולות שבין מקטעים סמוכים ומאמתת את הכיסוי לפני שהיא מסיקה את \textenglish{\texttt{Rails.Finite.finite\_all}}.

עבור $t\ge2485$, הקובץ \textenglish{\texttt{lean/Late.lean}} משייך את הקואורדינטה לאחת מחמש הרצועות שבסעיף~\ref{sec:five-bands}. ההוכחה משתמשת בתעודה לחצי ישר ברוחב 106 מתוך \textenglish{\texttt{Rule110LeftCheckpoints}}, בגזרת האתר המחזורי \textenglish{\texttt{EtherWedge}}, בהוכחה שב־\textenglish{\texttt{Crystal50WedgeAll}}, ברצועה הצרה \textenglish{\texttt{LeftCrystal50}} וברצועת השפה הימנית ברוחב 20 מתוך \textenglish{\texttt{Rule110Right20}}. הקובץ \textenglish{\texttt{Crystal50WedgeAll.lean}} מוכיח כי התווך הגבישי והחלון בן 16 הסיביות נשמרים יחד עבור כל אחד משני ערכי הסיבית החיצונית הבאה. בניית האינווריאנט המשותף הזה כוללת אימות של טבלאות המעבר הסופיות ושל תנאי ההתאמה בתפר. כיסוי חמשת המקרים נותן את \textenglish{\texttt{Rails.late\_all}}.

\subsection{היקף האימות הממוחשב}
לצורך המונוגרפיה בנינו מחדש את מודול השורש ואת המודולים שהוא מייבא על \textenglish{macOS arm64}, באמצעות גרסת \textenglish{Lean 4.35.0} שנקבעה בקובץ \textenglish{\texttt{lean-toolchain}}. \src{verification/monograph/geometry-axioms.log}{יומן אימות הגאומטריה} מתעד עיבוד של 263 מודולים וכיסוי של כל $5\,147\,263$ הזוגות בתחום הסופי, ב־100 מקטעים. הן עבור \textenglish{\texttt{Rails.late\_all}} והן עבור \textenglish{\texttt{Rails.essential\_iff}}, הליבה דיווחה על \textenglish{\texttt{[propext, Classical.choice, Quot.sound]}}. \src{verification/v2/final.log}{יומן הגרסה הקודמת} מתעד את אותן רשימות אקסיומות. טענה~\ref{prop:ring_cert} קובעת בנפרד קריטריון סופי לבדיקת מהותיות מתמדת בטבעת סופית כלשהי. הערה~\ref{rem:ring5} נותנת דוגמה נגדית להתמדה בטבעת של חמישה תאים. המניינים עבור הטבעת של 20 תאים מתקבלים ממנייה סופית ממצה.

הפקודה \textenglish{\texttt{make verify}} מחוללת מחדש את התעודה הסופית, בודקת את כלל הייבואים הישירים והעקיפים של \textenglish{\texttt{Rule110Exact}} לפי סדר התלויות, ומאמתת את דוח האקסיומות של משפט השורש. הפקודה \textenglish{\texttt{make verify-series}} בודקת גם את מודול הצופה ואת מרשמי הטענות והעדויות.

\subsection{אריחים אוטונומיים חד־צדדיים ושתי רצועות שפה שאינן חסומות בזמן}\label{app:half-ray-tiles}

החלה חוזרת של למה~\ref{lem:three_step} מראה שקבוצת התלות של הפלט בתא $s$ לאחר $3m$ צעדים מקיימת $\mathrm{Dep}_{3m}(s) \subseteq [s - 2m, s + 3m]$ לכל התצורות.

\begin{lemma}[סגירות על חצי ישר]\label{lem:one_sided}
יהי $A_t = F^t(A_0)$, יהי $p > 0$, ונניח כי $a(t + p) = a(t) + \delta$. נניח שבכל זמן בסיס $t_0 \le t < t_0 + p$ מתקיים $A_{t+p}(s + \delta) = A_t(s)$ לכל $s \ge a(t)$. נניח גם שהצמצום של $A_{t+p}$ אל $[a(t+p), \infty)$ תלוי רק בצמצום של $A_t$ אל $[a(t), \infty)$ בכל צעד באורך $p$. אזי $A_{t+p}(s + \delta) = A_t(s)$ לכל $t \ge t_0$ ולכל $s \ge a(t)$. הטענה המקבילה נכונה לחצאי ישר שמאליים $s \le a(t)$.
\end{lemma}

נחיל את למה~\ref{lem:one_sided} על החזית הפיזית הימנית $e(t) = L(t)$ עם $(\Delta t, \Delta s) = (30, +20)$. כאן $e(t+30) = e(t) + 20$, וקבוצת התלות של האריח החדש שרוחבו $W$ כלולה ב־$[e(t) - W + 1, e(t) + 50]$. הקצה השמאלי של הקטע מחריג את התחום הפנימי שאינו ידוע, $s < e(t) - W + 1$. החלת אותו נימוק על החזית השמאלית $s = -t$ עם $(\Delta t, \Delta s) = (40,-40)$ נותנת את התוצאה הבאה לזמן בלתי חסום.

\begin{theorem}[שתי רצועות שפה]\label{thm:dual_strips}
לכל זמן שלם $t \ge 300$,
\begin{equation}\label{eq:dual_strips}
[-L(t), -L(t) + 105] \cap \mathbb{Z} \;\cup\; [t - 7, t] \cap \mathbb{Z} \;\subseteq\; D_t.
\end{equation}
\end{theorem}

התעודה לרצועה השמאלית, שרוחבה $106$, משתמשת ב־$20$ מופעים של הרקעים $(0001111001)$, $(1011)$ ו־$(001011)$, בעומק אריח של $167$ תאים ולאורך $30$ זמני הבסיס $t = 300, \dots, 329$. התעודה לרצועה הימנית, שרוחבה $8$, משתמשת בשישה מופעים עם $(\Delta t, \Delta s) = (40,-40)$ לאורך זמני הבסיס $t = 300, \dots, 339$. \emph{חישוב הנושא הוכחה} על וקטורי סיביות מטיפוס \textenglish{\texttt{BitVec 1320}} מאמת את מצבי הביניים. הליבה של \textenglish{Lean 4.35.0} בודקת כל מעבר $F_{1320}^{h_i}(z_i) = z_{i+1}$ עבור $h_i \in \{1, 10\}$.

\subsection{קרניים בליסטיות ומגבלותיהן של תעודות מקומיות}

דרך נוספת להרחיב את הכיסוי לזמנים גדולים כרצוננו היא לשמר הבדל לאורך מסלול בדיד במרחב ובזמן.

\begin{lemma}[התמדה על חצי ישר ימני]\label{lem:ballistic_ray}
נקבע את שפת חצי הישר $s_0\in\mathbb Z$. נניח שמסלול $A$ מקיים $\forall s \ge s_0,\; A_{T+3}(s + 2) = A_T(s)$ בזמן $T$. אזי $\forall q \ge 0,\; \forall s \ge s_0,\; A_{T+3q}(s + 2q) = A_T(s)$. בפרט, אם גם $B$ וגם $\widetilde B$ מקיימים את התנאי הזה ואם $B_T(s_0) \ne \widetilde B_T(s_0)$, אז $B_{T+3q}(s_0 + 2q) \ne \widetilde B_{T+3q}(s_0 + 2q)$ לכל $q \ge 0$.
\end{lemma}

נתבונן בשתי המשפחות החשבוניות של נקודות ההתחלה על הרקעים $(0000101111)$ ו־$(001011)$, עבור $0 \le n \le 32$:
\begin{equation}
T_A(n) = 594 + 882n,\quad s_A(n) = 256 + 378n; \qquad
T_B(n) = 884 + 882n,\quad s_B(n) = 380 + 378n,
\end{equation}
בהתאמה. החישוב אימת את $66$ נקודות ההתחלה האלה. מכל אחת יוצאת קרן בליסטית אינסופית של קואורדינטות מהותיות בעומק קבוע ביחס לחזית $e(t) = L(t)$. העומקים הם $d_A(n) = 140 + 210n$ ו־$d_B(n) = 210 + 210n$, בהתאמה.

המודולים הבאים, \textenglish{\texttt{StaircaseA}} ו־\textenglish{\texttt{StaircaseB}}, מוכיחים את הנוסחאות לשתי המשפחות לכל $n,q\ge0$ באמצעות חיבור של שלוש שורות תואמות מקומית למסלול מדויק. התווכים המחזוריים בחזיתות של משפחות \textenglish{A} ו־\textenglish{B} מכסים $45\,255$ מתוך $66\,150$ מחלקות השאריות. הוספת משפחה שלישית מעלה את המניין ל־$53\,550$. המניינים האלה מתייחסים לטבלת השאריות בגודל $315\times210$. החלת רשומה מן הטבלה על קואורדינטה פיזית מחייבת גם לוודא שהנקודה נמצאת בגזרה המתאימה.

דוגמאות נגדיות מצביעות על שלוש מגבלות של התעודות האלה.
\begin{itemize}
    \item \textbf{אריחים מקומיים אינם מספיקים לצעד האינדוקציה מ־$n$ ל־$n+1$.} עבור הצעד $(\Delta T, \Delta s) = (882, +378)$, חסם התלות האוניברסלי הוא $\mathrm{Dep}_{882}(s) \subseteq [s - 588, s + 882]$. במטרה המוזזת מתקבל אפוא הקטע $[s - 210, s + 1260]$. הדוגמאות הנגדיות בקובץ \textenglish{\texttt{check\_macro\_boundary.py}} מקבעות אריחים ברדיוס $64$ (ב־$T=594$) או $186$ (ב־$T=884$). בכל אחד מהמקרים, שינוי תאים מחוץ לאריח משנה את סיבית המטרה לאחר $882$ צעדים.
    \item \textbf{בגזרה הנצפית במשך זמן סופי עשויים להתגלות בהמשך פערים.} ברקע $(0001111001)$, הגזרה השמאלית שמקיפה את המהירויות ב־$[-0.66, -0.64]$ נראית מכוסה עד $t = 2000$. הרצה ממושכת יותר מגלה פער.
    \item \textbf{הישנות עד כדי הזזה, כשלעצמה, אינה מוכיחה כיסוי.} אפילו הישנות עד עומק של $261$ תאים עבור $(0000001001)$ אינה מבטיחה כיסוי של כל מחלקות השאריות. הוכחת הכיסוי משלבת כמה מופעים.
\end{itemize}

\subsection{השלמה קונסטרוקטיבית של הגזרה השמאלית הבלתי חסומה (\textenglish{\texttt{wedge\_all}})}

כדי להביא בחשבון את תנועותיהם השונות של השפה השמאלית של חרוט האור (שמהירותה $-1$), של הגביש שמחזורו $50$, ושל חזית ההפרעה, נאנדקס את הזמנים הזוגיים באמצעות $n \in \mathbb{N}$ ונגדיר
\begin{equation}
\mathrm{tm}(n) = 2450 + 2n, \qquad a(n) = -350 - n.
\end{equation}
השפה הימנית $a(n)$ נעה תא אחד שמאלה בכל שני צעדים, במהירות $-1/2$. לכל אחד מששת המופעים $p \in \{\mathrm{p0}, \mathrm{p1}, \mathrm{p2}, \mathrm{p3}, \mathrm{p7}, \mathrm{p9}\}$, נכתוב $\widetilde B_p(t,s)=\widetilde B^p_t(s)$. לאינווריאנט בשכבה $n$ שני חלקים. המסלול מסכים עם גביש הייחוס $\mathrm{Cm}_p$ בקטע שמשמאל ל־$a(n)$, וחוצץ השפה שלו, בן $16$ סיביות, שווה למילה המאוחסנת בצומת של גרף מעברים סופי:
\begin{align}
\mathrm{Medium}(p,n)\ \iff\ {}&
\forall s\in[-\mathrm{tm}(n)+1300,a(n))\cap\mathbb{Z},\notag\\
&\widetilde B_p(\mathrm{tm}(n),s)=\mathrm{Cm}_p(\mathrm{tm}(n),s), \\
\mathrm{Window}(p,n,i)\ \iff\ {}&
\forall k\in\{0,\dots,15\},\notag\\
&\widetilde B_p(\mathrm{tm}(n),a(n)+k)=\mathrm{testBit}(\mathrm{word}(p,n,i),k).
\end{align}

\begin{lemma}[מתמר דטרמיניסטי בן 16 סיביות והממשק שלו עם הרצועה השמאלית]\label{lem:transducer}
לכל מופע $p$, הטבלה \textenglish{\texttt{Crystal50Packed}} מגדירה גרף ששכבותיו חוזרות במחזור $50$. כל רשומה תופסת $36$ סיביות. $16$ הסיביות הפחות משמעותיות מקודדות את החלון $\mathrm{word}(p, n, i)$, והסיביות $16\dots 25$ ו־$26\dots 35$ מקודדות את המעברים לשני הערכים האפשריים של סיבית הקלט מימין, $r \in \{0, 1\}$.
\begin{enumerate}[label=(\alph*)]
    \item בצעד $n$, ההקשר השמאלי בן $4$ הסיביות $\mathrm{ctx}(p, n \bmod 50)$ תופס את המיקומים $a(n) - 4 + m$ עבור $m \in \{0, 1, 2, 3\}$. עומקם הוא $t + s = 2096 + n + m$, והם נמצאים בקטע שכבר מכוסה על ידי $\mathrm{Medium}(p, n)$.
    \item החלה כפולה של פעולת הצמצום הסיבית, $\mathrm{shrink}(\mathrm{shrink}(\mathrm{input}(p, n, i, r)))$, מחשבת את $17$ סיביות התצורה בקטע $[a(n)-2, a(n)+14]$ בזמן $\mathrm{tm}(n+1)$. הסיבית הפחות משמעותית (עבור $k=0$) נותנת את $\mathrm{Cm}_p(\mathrm{tm}(n+1), a(n) - 2)$. יתר $16$ הסיביות ($k=1\dots 16$) מרכיבות את המילה בצומת $z$ של השכבה הבאה, שהחלון שלה מתחיל ב־$a(n+1) = a(n) - 1$.
    \item בשפה השמאלית של הרצועה, עומקי התאים $-\mathrm{tm}(n+1) + 1300 \le s < -\mathrm{tm}(n) + 1302$ הם $k = s + \mathrm{tm}(n+1) \in \{1300,1301,1302,1303\}$. המשפט \textenglish{\texttt{seam}} מתוך \textenglish{\texttt{Crystal50Seam}} מכסה תאים אלה משום ש־$t + s$ נשמר תחת ההעתקה $(t, s) \mapsto (t + 50, s - 50)$.
\end{enumerate}
\end{lemma}

\begin{theorem}[\textenglish{\texttt{Rails.Crystal50Wedge.wedge\_all}}]\label{thm:wedge_all}
לכל זמן שלם $t \ge 2485$ ולכל $s$ שלם המקיים
\begin{equation}\label{eq:wedge_region}
-t + 1302 \le s \le 59 - \lfloor t/2 \rfloor,
\end{equation}
מתקיים הפרדיקט \textenglish{\texttt{EssentialAt}}\,$t\,(-s)$. לזמן $t$ אין חסם עליון.
\end{theorem}

\begin{proof}[מתווה ההוכחה]
נגדיר $n = \lfloor (t - 2450) / 2 \rfloor$. עבור $t \ge 2485$ מתקיים $n \ge 17$, והשפה הימנית היא $59 - \lfloor t/2 \rfloor = -1166 - n = a(n) - 816$. בזמנים זוגיים $t = \mathrm{tm}(n)$, האינווריאנט $\mathrm{Medium}(p, n)$ נותן ישירות את השוויון $\widetilde B_p(t, s) = \mathrm{Cm}_p(t, s)$ לכל ששת המופעים $p$. בזמנים אי־זוגיים $t = \mathrm{tm}(n) + 1$, קטע התלות לצעד יחיד $[s - 1, s + 1]$ נמצא כולו בתוך $[-\mathrm{tm}(n) + 1300, a(n))$. לכן למת המקומיות \textenglish{\texttt{carry}} משמרת את ההתאמה גם בצעד זה. החלת \textenglish{\texttt{essential\_of\_cm}} הופכת את ההתאמה לעדות המהותיות הנדרשת.
\end{proof}

\subsection{כיסוי מלא על הישר האינסופי}\label{sec:five-bands}
נשתמש בקואורדינטת הקלט $j=-s$ ונכתוב $b=\lfloor t/2\rfloor$. עבור $t\ge2485$, חמש התוצאות הבאות חלות בקטעים עוקבים. שני הקצוות כלולים בכל שורה.
\begin{center}
\begin{tabular}{lll}
\toprule
קואורדינטה ראשונה & קואורדינטה אחרונה & תוצאה \\
\midrule
$-L(t)$ & $-L(t)+105$ & תעודה לחצי ישר ברוחב 106 \\
$-L(t)+106$ & $b-60$ & \textenglish{\texttt{EtherWedge.half\_sector}} \\
$b-59$ & $t-1302$ & \textenglish{\texttt{Crystal50Wedge.wedge\_all}} \\
$t-1301$ & $t-20$ & \textenglish{\texttt{LeftCrystal50.left\_strip}} \\
$t-19$ & $t$ & \textenglish{\texttt{right\_boundary\_width20}} \\
\bottomrule
\end{tabular}
\end{center}
שלושה מופעים של האתר שמחזורו 14 מבססים את השורה השנייה לכל $t\ge210$. השורה האחרונה מרחיבה את הרצועה הימנית הקודמת, שרוחבה שמונה. ב־$t=2485$, השורה השלישית כוללת רק את התא $j=1183$. כל שורה מתחילה בקואורדינטה העוקבת לזו שבה מסתיימת קודמתה. הלמה \textenglish{\texttt{Rails.late\_all}} מחברת את חמש ההוכחות.

הבדיקות המספריות הראשונות כיסו את $0\le t\le2000$. ניסוי מאוחר יותר, שהשתמש במופעים קבועים של שבע מילים, הגיע ל־$t=5000$. לצורך ההוכחה דרוש רק התחום הסופי הגובל בחמש הרצועות.

\paragraph{התעודה הסופית.}
עבור $0\le t<2485$, נשתמש בכל המופעים המחזוריים של שבע מילים. הזוגות שלהלן מציינים את האורך ואת הקידוד כמספר שלם של כל מילה, כאשר התא הראשון נמצא בסיבית הפחות משמעותית:
\[
(8,244),\ (10,632),\ (10,976),\ (10,576),\ (4,13),\ (6,52),\ (6,48).
\]
אנו מייצגים את הרקעים באמצעות וקטורי סיביות באורך 5520 והופכים סיבית יחידה במיקום 2760. המחזורים של כל שבע המילים מחלקים את 5520. בכל צעד זמן אנו מחשבים את האיחוד הסיבי של ההבדלים על פני כל הזוגות, ומוודאים שהוא מכיל את המסכה של הקטע הפיזי $[-t,L(t)]$. כך מכוסים $\sum_{t=0}^{2484}(2t+1-\lfloor t/3\rfloor)=5\,147\,263$ זוגות מרחב־זמן. עבור $s$ בקטע זה מתקיים $-5520<s-t$ וכן $s+t<5520$. לכן למת המקומיות \textenglish{\texttt{local\_obs}} מעבירה את התוצאה להיפוך סיבית יחידה על הישר האינסופי.

שבע המילים קובעות את השורה ההתחלתית. המחולל מספק תמונות מצב מועמדות בכל 25 צעדים. \textenglish{Lean} בודקת את מסכות הכיסוי בכל מקטע, ובכל מקום שבו יש מקטע עוקב היא מוודאת שמצב הסיום של המקטע שווה לתמונת המצב בתחילת המקטע הבא. במקטע האחרון יש עשר שורות. אינדוקציה על אורך המקטע משייכת לכל סיבית מכוסה במסכה זוג מסלולים מסוים. בדרך זו \textenglish{\texttt{Rails.Finite.finite\_all}} מוכיחה מהותיות לכל קואורדינטת קלט בתחום הסופי.

\begin{proof}[השלמת הוכחת משפט~\ref{thm:gap-free}]
משפט~\ref{thm:upper_bound} נותן את $D_t\subseteq I_t$. להכלה ההפוכה, נחיל את התעודה הסופית עבור $t<2485$ ואת חמש הרצועות עבור $t\ge2485$. הזזת כל זוג שנבדל בקואורדינטה הפיזית $s$ נותנת עדות לכך שקואורדינטת הקלט $-s$ מהותית.
\end{proof}

\subsection{מגאומטריה לתצפית חלקית}
קובץ המקור \textenglish{\texttt{observer/Rule110Observer.lean}} מייבא את \textenglish{\texttt{Rule110Exact}} ואת \textenglish{\texttt{Rule110FrontGlue}}. ההוכחה של \textenglish{\texttt{essential\_iff\_inCone}} מחילה את משפט השורש על ההעתק $j-s$, ואז מזיזה את התצורה כדי לקבל עדות עבור תא המטרה $s$. לכן חלון התלות המדויק של המטרה הוא
\[
  [s-L(t),s+t]\cap\mathbb Z.
\]
ההגדרה \textenglish{\texttt{CanPredict}} שואלת אם סדרת קריאות מאתרים שנקבעו מראש ועד למועד קבוע מספיקה לחישוב סיבית המטרה לכל קלט קביל. הלמה \textenglish{\texttt{canPredict\_iff}} מאפיינת את תנאי המספיקות: התשובה קבועה בכל קבוצה של קלטים בעלי אותה היסטוריית תצפית. הקובץ גם מספק דוגמאות נגדיות לטענה שכיסוי גאומטרי לבדו מספיק. נקודה עשויה להימצא בתוך החרוט הסיבתי גם כאשר קריאותיו של הצופה הזה אינן קובעות את ערך המטרה.

קוד המקור של הצופה אומת באמצעות \textenglish{Lean 4.35.0} עם שני הייבואים הקנוניים. \src{verification/monograph/observer-axioms.log}{יומן הצופה} מכיל את כל 11 התשובות מתוך 11 לפקודה \textenglish{\texttt{\#print axioms}}, החל ב־\textenglish{\texttt{canPredict\_iff}} וכלה ב־\textenglish{\texttt{coverage\_does\_not\_suffice}}. עבור כל משפט דיווחה הליבה במדויק על \textenglish{\texttt{[propext, Classical.choice, Quot.sound]}}. המניינים של \textenglish{OLMo} באים מהניסויים הנפרדים שבסעיף~\ref{sec:olmo}.

\section{מפרט שערי הבדיקה ואימות המדידות}
\label{app:gates}

בסעיף~\ref{sec:olmo} מובאות התוצאות המספריות למהדורה זו. התנאים שלהלן קובעים מתי מדידה תומכת בטענה המתאימה. כל השוואה מוגדרת עבור נקודת ביקורת אחת של המודל, מחלק אחד לאסימונים וממשק קבוע אחד. המפרט המדויק למהדורה זו מופיע ב\src{observer/gate_contract_ru.md}{מפרט הסדרה השנייה}.

\paragraph{מערכי הבדיקה והתקציב.}
נכתוב את שמונה סיביות הקלט בצורה $\mathbf{x}=LwR$. נמיין את המחרוזות האמצעיות $w\in\{0,1\}^6$ בסדר עולה של ערך הגיבוב \textenglish{SHA-256} של מחרוזת ה־\textenglish{ASCII} \verb!observer-v2|!, ואחריה שש הסיביות, לרבות אפסים מובילים. הדירוג מתחיל ב־1: המקומות \textenglish{1--8} מיועדים לפיתוח (\textenglish{dev}), \textenglish{9--20} לבדיקה משותפת של \textenglish{G1} ו־\textenglish{G2}, \textenglish{21--32} ל־\textenglish{G3} ו־\textenglish{33--44} ל־\textenglish{G4}. המקומות \textenglish{45--64} נשמרים בעתודה. לכל מחרוזת אמצעית נציג את המכפלה הקרטזית של ארבעת צירופי הקצוות וארבע ההיסטוריות, $(L,R,h,e)\in\{0,1\}^4$, בסדר לקסיקוגרפי. מתקבלות 16 שורות בכל בלוק, 128 שורות פיתוח ו־192 שורות בכל אחד ממערכי הבדיקה בני 12 הבלוקים שהוגדרו. מערך הפיתוח משמש לבחירת הניסוח, ההדגמות, השכבה וגודל ההתערבות. מחרוזות ההדגמה נשמרות בנפרד ממחרוזות ההערכה. יש להקפיא את המימושים המושווים לפני פתיחת מערך הבדיקה המשותף שנשמר בצד. שיטה שנבחרה על סמך תוצאות ממערך שכבר נפתח מחייבת מערך חדש.

\paragraph{קבלת שתי ההיסטוריות.}
תחילה נמדוד על פני אוצר האסימונים כולו את ההסתברויות להשלמות בנות אסימון יחיד \textenglish{\texttt{[space]A}} ו־\textenglish{\texttt{[space]B}}, בטמפרטורה 1. תנאי הסף הוא $P(A),P(B)\ge 0.01$. נשתמש ב־\textenglish{\texttt{random.Random(110)}} כדי לדגום בדיוק 1024 דגימות עם החזרה מההתפלגות הבינרית המותנית $P(A\mid A\text{ or }B)=P(A)/(P(A)+P(B))$. נרשום את שתי ההסתברויות המקוריות, את מלוא רצף הדגימות ואת מיקומה של ההופעה הראשונה של כל אחת מהאותיות \textenglish{A} ו־\textenglish{B}. בכל היסטוריה המודל מעבד את האסימון הראשון מן הסוג המתאים שנדגם בפועל. תמונת המצב נשמרת אחרי עיבוד האסימון ולפני קבלת המשוב. אם אחת ההסתברויות יורדת מתחת לסף או אם אחת האותיות אינה מופיעה, הקידומת מקבלת את המעמד \textenglish{\texttt{NOT\_EVALUABLE}}. אילוץ של אות מסוימת או דגימת מקבץ נוסף היו משנים את תנאי הניסוי.

\paragraph{קריטריונים משותפים.}
בכל שורה המחייבת מדידה מספרית, לרבות בקרות עם מצב נייטרלי, שוחזר או הוחלף, יש לרשום את ערך היעד, את האסימון בעל הערך המרבי באוצר האסימונים כולו, את $\ell_1-\ell_0$ ואת $q=P(0)+P(1)$. ההסתברויות מחושבות באמצעות נרמול על פני אוצר האסימונים כולו, בדיוק \textenglish{FP32} או גבוה ממנו. בכל שורה נדרשים ערכים סופיים וכן $q\ge 0.80$. כאשר ההפרש חיובי ממש או שלילי ממש, ההסתברות של הסיבית המנצחת היא $>0.40$, וההסתברות של כל אסימון מחוץ לזוג היא לכל היותר $0.20$. הפרש אפס נחשב לכישלון, והשורה נותרת במכנה. בלוק נחשב שלם רק כאשר כל 16 מתוך 16 הפלטים נכונים. ספי ההצלחה הם 6 מתוך 8 בלוקים שלמים במערך הפיתוח ו־9 מתוך 12 במערך הבדיקה שנשמר בצד. במקרים שבהם $L\ne R$ יש לדווח על הזוגות שנפתרו בנפרד לכל ערך של $e$. השיעור המצוין, 75\%, מתייחס לבלוקים שלמים. שימוש חוזר בקידומת אינו יוצר דגימות בלתי תלויות של מופעי המודל.

\paragraph{\textenglish{G0} ו־\textenglish{G0R}.}
בדיקת \textenglish{G0} מאמתת את ערכי הגיבוב של המשקולות ומחלק האסימונים, את הייצוג באסימון יחיד של \textenglish{A/B/X} ושל \textenglish{0/1} בהקשרים בפועל, את ההרכב המדויק של מזהי האסימונים, את צורת הטנזורים, את טיפוס הנתונים ואת מיקומי \textenglish{RoPE}. בדיקות האבחון נערכות על עותקים לשימוש חד־פעמי. קריאת הבחירה הקודמת נעצרת לפני שהתשובה נכתבת חזרה. אותה שאילתה ראשית חייבת להחזיר לוגיטים זהים על פני אוצר האסימונים כולו, עם בדיקת האבחון ובלעדיה. בתמונות המצב המקוריות של שני הענפים נשמרים הבתים של כל רשומות \textenglish{K/V} הפעילות, ההיסט וסדר המיקומים הנראים, לרבות אינדקסים בחוצץ המעגלי. ההשוואה חלה על הנתונים הלוגיים במיקומים הפעילים. קיבולת זיכרון מוקצית שאינה בשימוש מחוץ לטווח זה אינה נכללת בהשוואה.

לפני הבחירה, \textenglish{G0R} מחליפה באופן הפיך את מצב מיקום הבחירה בכל השכבות, ללא תלות בציון ההתנהגותי. בהינתן היסטוריה קודמת משותפת $U$, תמונות המצב הן $C_A=(U,a)$ ו־$C_B=(U,b)$. לכל נוסח משוב ולכל מחרוזת חדשה, נבדוק הן את המסלול $A\to B\to A$ והן את המסלול ההפוך המתחיל ב־\textenglish{B}. נבצע את ההחלפה ההפוכה על תוצאת ההחלפה הקדמית. את הקריאות יש לבצע על עותקים מן התוצאה. נשווה את מצב הביניים להיסטוריית התורם שלא שונתה, ואת מצב הסיום להיסטוריית הפתיחה. נשווה מטא־נתונים, רשומות \textenglish{K/V} פעילות ולוגיטים באותה פריסה. הפלטים חייבים להיות סופיים, בעלי צורות תואמות וזהים ברמת הבתים בטיפוס הנתונים המקורי. אחרי המרת הלוגיטים ל־\textenglish{FP32}, ההפרש המוחלט המרבי חייב להיות אפס. נבדוק את תמונות המצב להשוואה לפני הפעולה ואחריה. הרצת \textenglish{Base} הראשונה השתמשה בקידומת המקורית הקבועה \textenglish{\texttt{prefix(8)}}. לפני ההחלפה יש לבדוק שהקידומת המשותפת אכן זהה ושקיים לפחות בית פעיל אחד ששונה במיקום הבחירה. אם אי אפשר להבחין בין הענפים כבר במיקום זה, \textenglish{G0R} מחזירה \textenglish{\texttt{NOT\_EVALUABLE}} לקידומת. במערך הפיתוח, שתי השוואות לוגיטים בכל אחת מ־128 השורות נותנות 256 בדיקות.

\paragraph{\textenglish{G1} ו־\textenglish{G2}.}
בדיקת \textenglish{G1} בוחנת חישוב ממשוב טקסטואלי בהיסטוריות שלא שונו. נוסף על המדידות המספריות המחייבות, יש להשתמש בעותקים נפרדים עבור \textenglish{D1} (להיזכר ב־$h$ לפני המשוב), \textenglish{D2} (לנקוב בכלל המתוקן עבור ארבעת הזוגות $h,e$ ללא המחרוזת החדשה) ו־\textenglish{D3} (להחיל כלל ששמו ניתן במפורש באפיזודה חדשה). יש לדווח על בדיקות האבחון שפלטן אות בנפרד מן המשימה המספרית הראשית. הצלחה בבקרה הטקסטואלית מחייבת עמידה בפורמט ובשיעור הבלוקים השלמים שנקבעו. ב־\textenglish{G2}, לאחר \textenglish{G0}, \textenglish{G0R} ו־\textenglish{G1}, יש לחזור על השוואות התורם והשחזור במערך \textenglish{9--20}. נוסיף מטמון משותף ובו אסימון נייטרלי \textenglish{X} באותו מיקום ובאותו אורך אסימון. הזהות עם התורם שלא שונה נובעת מן השוויון $P_B(U,a)=(U,b)$. לפיכך, תוצאת הבדיקה הטכנית החיובית היא \textenglish{\texttt{REPLAY\_OK}}. ההתנהגות בהיסטוריות שלא שונו כבר נמדדת ב־\textenglish{G1}.

\paragraph{\textenglish{G3}.}
לאחר \textenglish{G1} וקבלת \textenglish{\texttt{REPLAY\_OK}}, נחליף את האות הטקסטואלי בהתערבות יחידה באסימון נייטרלי לפני כתיבת רשומות ה־\textenglish{K/V} שלו. לפני הבדיקה נפרט את האסימון, המיקום, השכבה, נקודת ההתערבות, מקור הווקטור וגודלו. נייצג וקטור יחיד, סופי ושאינו אפס, $v$, בטיפוס הנתונים של הפעולה. היפוך סימן מדויק נותן את הפעימות $v_0=v$, $v_1=-v$. עבור ערך נתון של $e$, שתי ההיסטוריות מקבלות טנזורים זהים ברמת הבתים. נבדוק את הצורה, את הטיפוס, את היפוך הסימן לכל רכיב ואת השוויון בנורמות האוקלידיות שאינן אפס, המחושבות ב־\textenglish{FP64} באותו סדר. לכל ענף בנפרד נרשום את נורמת התוספת אחרי העיגול בפועל, $\operatorname{round}(r+v_e)-r$. נכבה את ההתערבות לאחר חיבור יחיד זה. כאשר ההתערבות היא אפס, שני ענפי $e$ בערכים קבועים של $h,\mathbf{x}$ חייבים לקבל קלטים ולוגיטים זהים. נדווח על דיוקם ועל מסת ההסתברות שלהם. תוצאה חיובית עם הפעימה מחייבת את הפורמט המשותף ולפחות 9 מתוך 12 בלוקים שלמים במערך \textenglish{G3}.

\paragraph{\textenglish{G4}.}
לפני הצגת $\mathbf{x}$, ארבעת התורמים מעבדים אותו רצף נייטרלי לאחר $h,e$. לתבנית חדשה נעביר רק את רשומות ה־\textenglish{K/V} במיקומים ובשכבות שנקבעו מראש בתוך רצף זה, ונסמן אותן $M_{he}$. כל יתר המיקומים הפעילים, האסימונים, האורכים והמטא־נתונים זהים אצל ארבעת הנמענים. לנמענים אין גישה ישירה ל־$h,e$ או לתוצרים נגזרים אחרים של היסטוריית התורם. את המצב הנייטרלי $M_*$ נגדיר פעם אחת מתוך היסטוריה נייטרלית משותפת. הערוץ הטקסטואלי מחייב את \textenglish{G1} ואת \textenglish{G2}. ערוץ הפעימה מחייב גם את \textenglish{G3}. כל הקריטריונים הבאים חייבים להתקיים ב־12 המחרוזות האמצעיות:
\begin{enumerate}
\item נמען שמקבל את מלוא ההיסטוריה ואת אותו רצף ביניים עומד בקריטריון הפורמט ופותר לפחות 9 מתוך 12 בלוקים שלמים. אם הבקרה הזאת נכשלת, כישלון בהעברה אינו מכריע אם המידע הדרוש היה קיים אצל התורם.
\item לאחר הכנסת $M_*$, כל ארבעת הקלטים, רשומות ה־\textenglish{K/V} הפעילות והלוגיטים, עבור $\mathbf{x}$ קבוע, זהים ברמת הבתים. בתמונת מצב נגזרת אחת, החזרת $M_{he}$ השייך לה לאחר הכנסת מצבו של תורם אחר משחזרת את הלוגיטים המקוריים.
\item העברת $M_{he}$ של כל תורם עצמו עומדת בקריטריון הפורמט ובסף של 9 מתוך 12 בלוקים שלמים. בכל בלוק נכללים כל ארבעת התורמים וכל ארבעת צירופי הקצוות.
\item ב־96 השורות עם $L\ne R$, נסמן ב־$C_M$ וב־$C_*$ את מספר ההחלטות הנכונות לאחר העברה ולאחר מחיקה, בהתאמה. נדרשים פורמט תקין, $C_*\le48$, אפס מתוך 12 בלוקים שלמים לאחר המחיקה, וכן $(C_M-C_*)/96\ge 0.25$. נספור בנפרד שורות שתוקנו ושורות שנפגמו. כאשר נדרשת תשובה בינרית חד־משמעית אחת, הקלט המשותף מניב בדיוק 48 מתוך 96 תשובות נכונות. תשעה בלוקים שלמים במצב העברה מבטיחים לפחות 72 מתוך 96 בשורות שבהן הקצוות שונים.
\end{enumerate}

\paragraph{\textenglish{G4-1}.}
לפני שנראה $\mathbf{x}$ כלשהו, השולח בוחר בחופשיות אסימון אחד בפענוח חמדני על פני אוצר האסימונים כולו. מתקבלות רק שתי ההשלמות שנקבעו מראש, \textenglish{A/B}, שכל אחת מהן בת אסימון יחיד. שוויון במקום הראשון או בחירת אסימון אחר מניבים \textenglish{\texttt{FAIL\_FORMAT}}, ללא ניסיון חוזר. נרשום את ארבע ההודעות בסדר ההיסטוריות $00,01,10,11$. המפרט המקורי של \textenglish{Base} קובע את \textenglish{\texttt{prefix(8)}}, את המשוב המילולי \textenglish{correct/wrong} ובקשה לרשום את הכלל המתוקן. לתנאי \textenglish{Instruct} הדו־שלבי שנמדד יש הנחיות משלו וסף נוסף $P(A)+P(B)\ge0.80$, כפי שנקבע ב\src{observer/two-step-interface-04.json}{קובץ המפרט של חלופה \textenglish{A}}. יש להתייחס לתנאים של שתי נקודות הביקורת בנפרד.

כדי לשחזר את המפרט המקורי של \textenglish{Base}, לאחר בחירת השולח יש לתת לו \verb!\nFeedback: correct! או \verb!\nFeedback: wrong!, ואחריה \verb!\nRecord the corrected rule. Rule:!. כאן \verb!\n! מייצג מעבר שורה אחד. הקידומת החדשה של הנמען כוללת שתי שורות:
\begin{english}
\begin{verbatim}
Rule A selects the first bit. Rule B selects the last bit.
Rule:
\end{verbatim}
\end{english}
יש להכניס את המזהה הנבחר, הכולל רווח מוביל, מייד אחרי הנקודתיים. אחר כך יש להוסיף \verb!\nBits: !, את שמונה הסיביות מופרדות ברווח יחיד, ולבסוף \verb!\nAnswer: ! ובו הרווח המסיים.

רק המזהה הנבחר נמסר לנמען באפיזודה החדשה. הנמען מחשב מחדש את רשומות ה־\textenglish{K/V} עבור מזהה זה. יש לאמת את מזהי האסימונים של הקידומת המשותפת ושל מחרוזת השאילתה מול החלוקה לאסימונים של הטקסט השלם. ממשק החוץ אינו מעביר את רשומות ה־\textenglish{K/V} של השולח, את הלוגיטים שלו, את אינדקס הענף, את זמן התגובה או את פרמטרי הדגימה. בחירת האסימון הנושא את המסר או הענפים על סמך $h,e,p$, התשובות הנכונות או משתנים המשמשים להם תחליף אסורה. בבקרת המחיקה נציג את אותם קלטים בכל העולמות עם אסימון קבוע יחיד, \textenglish{A}, שנקבע מראש. תוצאת הבדיקה \textenglish{\texttt{ZERO\_ERROR\_ON\_PANEL}} דורשת ארבע הודעות קבילות, בדיקות לגבול הערוץ ולמחיקה, פורמט מספרי תקין ו־12 מתוך 12 בלוקים שלמים. כאשר \textenglish{G4} ו־\textenglish{G4-1} נבדקות יחד, יש להקפיא את שתי הגרסאות לפני הפתיחה הראשונה של מערך \textenglish{33--44}.

\paragraph{סדר קדימות של תוצאות הבדיקה ויכולת שחזור.}
שגיאה טכנית, ערך מספרי שאינו סופי או הפרה של דרישת זהות הבתים מניבים \textenglish{\texttt{INVALID}}. אחריהם באים, בסדר קדימות יורד, \textenglish{\texttt{NOT\_EVALUABLE}}, \textenglish{\texttt{FAIL\_FORMAT}} ו־\textenglish{\texttt{FAIL\_DEV}/\texttt{FAIL\_TEST}}. בדיקה התלויה בקודמת, ולא הגיעו לבצעה, מקבלת \textenglish{\texttt{NOT\_RUN}}. ממשק מתוקן או נקודת ביקורת שונה מגדירים תנאי נפרד.

\noindent\begin{minipage}{\linewidth}
יש להריץ את שתי הבדיקות הבאות מתיקיית השורש של החבילה המוקפאת.
\begin{english}
\begin{verbatim}
python3 -B observer/verify.py
python3 -B observer/verify_part_iii.py
\end{verbatim}
\end{english}
\end{minipage}\par
המאמת הראשון מחשב מחדש את הבניות הסופיות, בודק את קובץ המפרט של המהדורה שפורסמה עבור הקורפוס ההיסטורי, מאמת את כתובות המרשם ומשחזר את הטבלה מקובץ ה־\textenglish{SVG}. המאמת השני קורא את דוחות הפיתוח שנשמרו, מאמת את מערכי התנאים, ההסתברויות והלוגיטים שתועדו ומחשב מחדש את הדיוק, את הבלוקים השלמים, את זוגות ההיסטוריות המנוגדות, את תוצאות האבחון ואת השפעת המחיקה. עבור \textenglish{G0/G0R} הוא בודק מחדש את המטא־נתונים וערכי הגיבוב שתועדו להשוואות בזוגות של הלוגיטים בטיפוס הנתונים המקורי. הרצה חוזרת של ההיסק על \textenglish{GPU} והשוואה של הטנזורים החדשים בית אחר בית היו מהוות מדידה חדשה.

\section{מרשם הטענות, מקור העדויות ומשבצות פתוחות}
\label{app:registers}

ב\src{observer/hypotheses.tsv}{מרשם ההשערות} הקנוני יש 23 שורות. בכל שורה מופיעים מזהה, תלויות, נוסח מדויק של הטענה, הנחות, מעמד, הפניה לעדות ותנאי לבחינה מחדש. התלויות מופיעות לפני הטענות התלויות בהן. קבוצת ההיסטוריות הקבילות היא חלק מנוסח הטענה. לדוגמה, הכישלון בהכללת הדוגמה המקורית בת 11 הסיביות לכל המילים מתיישב עם הדיוק שלה בארבע המילים שנבחרו.

\begin{center}
\begin{tabularx}{\linewidth}{@{}l r X@{}}
\toprule
מעמד במרשם & מתוך 23 & מזהים \\
\midrule
הוכחה & 12 & \textenglish{OBS.01--03, R110.01, R110.05--09, R110.12, LOOP.01, E.01} \\
הפרכה & 4 & \textenglish{R110.02--04, R110.10} \\
אימות במסגרת הניסוי & 3 & \textenglish{AI.01--03} \\
פתוח & 2 & \textenglish{R110.11, AI.04} \\
\textenglish{\texttt{DECLARED\_NOT\_DERIVED}} & 1 & \textenglish{E.02} \\
\textenglish{\texttt{UN}} & 1 & \textenglish{E.03} \\
\bottomrule
\end{tabularx}
\end{center}

המעמדות \emph{הוכחה} ו־\emph{הפרכה} חלים על הטענה עם כל ההנחות המנויות בשורה שלה. העדויות בשורות השונות הן מסוגים שונים: הוכחה מתמטית, בנייה סופית הניתנת לאימות חישובי או מנייה ממצה. הטענות \textenglish{AI.01--03} מתעדות, בהתאמה, את הקשר בין תשובת העוזר לאות שהוצג לו, את ההבדל בין ערוץ חי לערוץ מוקלט בשאילתות חדשות ואת עדכונה של אסטרטגיה קודמת במשפחת המשימות שנבדקה. מקורן של שורות אלה בניסויים נפרדים בעוזר. לספירות של מודל \textenglish{OLMo} הקפוא בסעיף~\ref{sec:olmo} יש דוחות ומערכי בדיקה משלהן. \textenglish{E.01} מוכיחה שהפרדיקט הנוסף אינו נקבע באופן חד־משמעי. \textenglish{E.02} מתעדת הנחה נוספת, ו־\textenglish{E.03} מציינת שמעמדו של מופע מסוים עדיין פתוח.

\paragraph{36 כתובות לטענות.}
\src{observer/leaf/leaf.tsv}{טבלת המשבצות} מצרפת שישה תחומים, \textenglish{S} (מקור), \textenglish{I} (זהות), \textenglish{M} (משמעות), \textenglish{T} (מעבר), \textenglish{A} (פעולה) ו־\textenglish{O} (תצפית), לשישה תפקידים תיאוריים: מוצא, מצב, נשא, טענה מקומית, עדות וגבול. לכתובות בטבלה זו אנו משתמשים בתחילית \textenglish{\texttt{leaf:}}. לדוגמה, \src{observer/leaf/leaf.tsv\#L25}{\textenglish{\texttt{leaf:T6}}} מציינת את השורה בקובץ ה־\textenglish{TSV} עם \textenglish{\texttt{law=T6}}, המשבצת השישית בתחום המעבר. כתובת זו מפנה לתנאי של הערוץ הרועש (\textenglish{T3}). לעומת זאת, \hyperref[thm:conditional-bridge]{\textenglish{T6}} ללא תחילית מציין את המשפט המותנה שבסעיף~\ref{sec:boundary}. בגרסה זו 32 מתוך 36 השורות מסומנות \textenglish{\texttt{KNOWN}}. כל אחת מספקת ערך גיבוב \textenglish{SHA-256} של קובץ והפניה הניתנת לאימות להגדרה, לפונקציה או לקטע טקסט. ארבע מתוך 36 השורות הנותרות מסומנות \textenglish{\texttt{UNKNOWN}}, ושדות מקורות העדויות שלהן ריקים. השאלות המהותיות שבהן מתועדות כלהלן.
\begin{itemize}
\item \textenglish{\texttt{leaf:I6}} עוסקת בזהות המספרית של הסובייקט כאשר היסטוריה סיבתית מתפצלת. מוסכמת תיעוד הצאצאים מקצה מזהים לענפים.
\item \textenglish{\texttt{leaf:M6}} עוסקת בהצדקה בלתי תלויה לקשר הפסיכופיזי שנלקח כתנאי מספיק.
\item \textenglish{\texttt{leaf:A6}} דורשת לממש את ערוץ הצד כולו בתוך מערכת אחת של כלל 110 שהיא בינארית לחלוטין והומוגנית, עבור העדות המצוינת (\textenglish{R110.11}).
\item \textenglish{\texttt{leaf:O6}} עוסקת בערכו של הפרדיקט הפנומנלי עבור מופע מסוים ברגע המשוב (\textenglish{E.03}).
\end{itemize}
\src{observer/leaf/render_leaf.py}{בודק התרשים} מאמת את סדר הכתובות, את ערכי הגיבוב מול קובצי המקור, את קיום הסמלים, את ייצוג ה־\textenglish{TSV} הקנוני ואת האפשרות לשחזר אותו מתוך ה־\textenglish{SVG} בהתאמה מדויקת של כל הבתים. כך התרשים מספק תצוגה הניתנת לאימות של הטבלה.

\paragraph{מקור המהדורה ומעמדה.}
\src{SHA256SUMS}{מצאי המהדורה הראשי} כולל סכומי ביקורת של ה־\textenglish{PDF}, שני יומני \textenglish{Lean} וקובצי המקור. \src{observer/sha256.txt}{מצאי נפרד} מתעד את 48 הקבצים באוסף הצופה. נשמרה גם \src{observer/provenance/observer-original.sha256}{רשימת סכומי הביקורת ההיסטורית} בת 36 הרשומות מן המהדורה המקורית. חבילת הפרסום משתמשת ב־35 קבצים מקוריים מן הרשימה הזאת. הרשומה של מטלת המחקר הפנימית נשארת ברשימה לשם תיעוד המקור, אך המטלה עצמה הושמטה מהחבילה הציבורית. ערכי הגיבוב שפורסמו משקפים שינויים בניווט ובמיקומי העדויות. דוחות ה־\textenglish{JSON} השמורים מכילים את התצפיות המקוריות. חוזה הבדיקה שנעשה בו שימוש בהרצות קודמות נבדק מול הבתים של נוסחו במהדורה הקשורה לאותם דוחות.

\paragraph{תנאים להמשך העבודה.}
\textenglish{R110.11} דורשת קידוד הניתן לאימות של הקלט, האחסון, המסירה והמענה בתוך מערכת דינמית בינארית אחת, לצד מסלולים עבור כל ארבע ההיסטוריות. \textenglish{AI.04} דורשת משימה חדשה ובדיקת העברה של המצב הפנימי, כאשר עמדות המקור אינן נגישות, מלוא ההיסטוריה זמין בביקורת החיובית, ונכללות החלפה הפיכה ומחיקה. נספח~\ref{app:gates} מפרט את החוזה המספרי של \textenglish{G4}. כדי לאתר את ההעברה בתת־מרחב נתון, יש להגדיר לפני הניסוי את השכבה, את ההתערבות, את מדד התיווך, את החלופות שיוחזקו תחת בקרה ואת הסבולת. תוצאה עבור נשא טקסטואלי יחיד משאירה תנאי זה פתוח.

בדיקה מלאה של הייבואים הקנוניים של הצופה בוצעה בעת הכנת המונוגרפיה באמצעות \textenglish{Lean~4.35.0} (נספח~\ref{app:lean}). מעמדיהם של מערכי \textenglish{OLMo} מתייחסים למצב הניסוי בשלב הפיתוח. השערת הפלסטיות המתמשכת דורשת שינוי בארגון הנשא בין אפיזודות והדגמת השפעה סיבתית של אותו שינוי על עיבוד קלט בהמשך, כאשר ההיסטוריה החיצונית נתונה לבקרה. \textenglish{E.02} דורשת נימוקים בלתי תלויים לבחירת הנחת הגישור, או תחזיות אופרטיביות שונות לגשרים מתחרים בצירוף הנחות קישור מפורשות. \textenglish{E.03} נותרת במעמד \textenglish{\texttt{UN}} בגרסה זו.

\clearpage
\begingroup\small

\begin{thebibliography}{9}
\bibitem{cook2004}
\textenglish{M.~Cook, ``Universality in elementary cellular automata,'' \emph{Complex Systems}, vol.~15, no.~1, pp.~1--40, 2004.}

\bibitem{martinez2006}
\textenglish{G.~J.~Mart\'inez, H.~V.~McIntosh, and J.~C.~Seck-Tuoh-Mora, ``Gliders in Rule 110,'' \emph{International Journal of Unconventional Computing}, vol.~2, no.~1, pp.~1--49, 2006.}

\bibitem{shalizi2001}
\textenglish{C.~R.~Shalizi and J.~P.~Crutchfield, ``Computational mechanics: Pattern and prediction, structure and simplicity,'' \emph{Journal of Statistical Physics}, vol.~104, pp.~817--879, 2001.}

\bibitem{wolfram1986}
\textenglish{S.~Wolfram, ``Tables of cellular automaton properties,'' in \emph{Theory and Applications of Cellular Automata}, World Scientific, pp.~485--557, 1986.}

\bibitem{wolpert2008}
\textenglish{D.~H.~Wolpert, ``Physical limits of inference,'' \emph{Physica D: Nonlinear Phenomena}, vol.~237, no.~9, pp.~1257--1281, 2008.}

\bibitem{witsenhausen}
\textenglish{H.~S. Witsenhausen. The zero-error side information problem and chromatic numbers.
\emph{IEEE Transactions on Information Theory}, 22(5), 592--593, 1976.
\href{https://doi.org/10.1109/TIT.1976.1055607}{doi:10.1109/TIT.1976.1055607}.}

\bibitem{charpenay}
\textenglish{N.~Charpenay, M.~Le Treust, A.~Roumy. Side Information Design in Zero-Error Coding for Computing.
\emph{Entropy}, 26(4), 338, 2024.
\href{https://doi.org/10.3390/e26040338}{doi:10.3390/e26040338}.}

\bibitem{macar2026}
\textenglish{U.~Macar, L.~Yang, Z.~Wang, P.~Wallich, E.~Ameisen, J.~Lindsey.
\href{https://proceedings.mlr.press/v306/macar26a.html}{Mechanisms of Introspective Awareness}.
\emph{Proceedings of the 43rd International Conference on Machine Learning}, PMLR 306, 85403--85429, 2026.}

\bibitem{geiger2025}
\textenglish{A.~Geiger, D.~Ibeling, A.~Zur, M.~Chaudhary, S.~Chauhan, J.~Huang, A.~Arora, Z.~Wu, N.~Goodman, C.~Potts, T.~Icard.
\href{https://www.jmlr.org/papers/v26/23-0058.html}{Causal Abstraction: A Theoretical Foundation for Mechanistic Interpretability}.
\emph{Journal of Machine Learning Research}, 26(83), 1--64, 2025.}

\end{thebibliography}
\endgroup

\end{document}
